% Template for PLoS
% Version 3.8 Apr 2026
%
% % % % % % % % % % % % % % % % % % % % % %
%
% -- IMPORTANT NOTE
%
% This template contains comments intended 
% to minimize problems and delays during our production 
% process. Please follow the template instructions
% whenever possible.
%
% % % % % % % % % % % % % % % % % % % % % % % 
%
% Once your paper is accepted for publication, 
% PLEASE REMOVE ALL TRACKED CHANGES in this file 
% and leave only the final text of your manuscript. 
% PLOS recommends the use of latexdiff to track changes during review, as this will help to maintain a clean tex file.
% Visit https://www.ctan.org/pkg/latexdiff?lang=en for info.
%
%
% IMPORTANT
% Below are a few tips to help format manuscript according to our specifications. For more tips to help reduce the possibility of formatting errors during conversion, please see our LaTeX guidelines at http://journals.plos.org/plosone/s/latex
%
% There are no restrictions on package use within the LaTeX files except that no packages listed in the template may be deleted.
%
% Please do not include colors or graphics in the text.
%
% The manuscript LaTeX source should be contained within a single file (do not use \input, \externaldocument, or similar commands).
%
% % % % % % % % % % % % % % % % % % % % % % %
%
% -- FIGURES AND TABLES
%
% Please include tables/figure captions directly after the paragraph where they are first cited in the text.
%
% DO NOT INCLUDE GRAPHICS IN YOUR MANUSCRIPT
% - Figures should be uploaded separately from your manuscript file. 
% - Figures generated using LaTeX should be extracted and removed from the PDF before submission. 
% - Figures containing multiple panels/subfigures must be combined into one image file before submission.
% For figure citations, please use "Fig" instead of "Figure".
% See http://journals.plos.org/plosone/s/figures for PLOS figure guidelines.
%
% Tables should be cell-based and may not contain:
% - spacing/line breaks within cells to alter layout or alignment
% - do not nest tabular environments (no tabular environments within tabular environments)
% - no graphics or colored text (cell background color/shading OK)
% See http://journals.plos.org/plosone/s/tables for table guidelines.
%
% For tables that exceed the width of the text column, use the adjustwidth environment as illustrated in the example table in text below.
%
% % % % % % % % % % % % % % % % % % % % % % % %
%
% -- EQUATIONS, MATH SYMBOLS, OTHER SPECIAL FORMATTING
%
% - For bold symbols, use the \boldsymbol command and not \mathbf. 
% - For inline equations, please be sure to include all portions of an equation in the math environment.  For example, x$^2$ is incorrect; this should be formatted as $x^2$ (or $\mathrm{x}^2$ if the romanized font is desired).
% - Do not include text that is not math in the math environment. For example, CO2 should be written as CO\textsubscript{2} instead of CO$_2$.
% - Avoid capturing simple text as equations, for example $<$.
% - For inline equations, please do not include punctuation (commas, etc) within the math environment unless this is part of the equation.
% - When adding superscript or subscripts outside of brackets/braces, please group using {}.  For example, change "[U(D,E,\gamma)]^2" to "{[U(D,E,\gamma)]}^2". 
% - Do not use \cal for caligraphic font.  Instead, use \mathcal{}.
% - Avoid splitting for fragmenting a single equation into multiple objects/environments, for example $f(x)$ = $g(x)$ + $h(x)$.
% - Use math environment for all equations, do not mimic using text.
% - Avoid using single letters in math mode where possible.
% - Always use matching parentheses and delimiters, avoid mixing math and text brackets.
% - Insert spaces around inline equations to ensure proper display after conversion.
% - Do not use forced numbers for display equation labeling or suppress numbering with \nonumber. 
%
% % % % % % % % % % % % % % % % % % % % % % % % 
%
% Please contact customercare@plos.org with any questions.
%
% % % % % % % % % % % % % % % % % % % % % % % %

\documentclass[10pt,letterpaper]{article}
\usepackage[top=0.85in,left=2.75in,footskip=0.75in]{geometry}

% amsmath and amssymb packages, useful for mathematical formulas and symbols
\usepackage{amsmath,amssymb}

% Use adjustwidth environment to exceed column width (see example table in text)
\usepackage{changepage}

% textcomp package and marvosym package for additional characters
\usepackage{textcomp,marvosym}

% cite package, to clean up citations in the main text. Do not remove.
\usepackage{cite}

% Use nameref to cite supporting information files (see Supporting Information section for more info)
\usepackage{nameref,hyperref}

% line numbers
\usepackage[right]{lineno}

% ligatures disabled
\usepackage[nopatch=eqnum]{microtype}
\DisableLigatures[f]{encoding = *, family = * }

% color can be used to apply background shading to table cells only
\usepackage[table]{xcolor}

% array package and thick rules for tables
\usepackage{array}

% create "+" rule type for thick vertical lines
\newcolumntype{+}{!{\vrule width 2pt}}

% create \thickcline for thick horizontal lines of variable length
\newlength\savedwidth
\newcommand\thickcline[1]{%
  \noalign{\global\savedwidth\arrayrulewidth\global\arrayrulewidth 2pt}%
  \cline{#1}%
  \noalign{\vskip\arrayrulewidth}%
  \noalign{\global\arrayrulewidth\savedwidth}%
}

% \thickhline command for thick horizontal lines that span the table
\newcommand\thickhline{\noalign{\global\savedwidth\arrayrulewidth\global\arrayrulewidth 2pt}%
\hline
\noalign{\global\arrayrulewidth\savedwidth}}


% Remove comment for double spacing
%\usepackage{setspace} 
%\doublespacing

% Text layout
\raggedright
\setlength{\parindent}{0.5cm}
\textwidth 5.25in
\textheight 8.75in

% Bold the 'Figure #' in the caption and separate it from the title/caption with a period
% Captions will be left justified
\usepackage[aboveskip=1pt,labelfont=bf,labelsep=period,justification=raggedright,singlelinecheck=off]{caption}
\renewcommand{\figurename}{Fig}

% Please use the included `plos2025.bst` as your BibTeX style
\bibliographystyle{plos2025}

% Remove brackets from numbering in List of References
\makeatletter
\renewcommand{\@biblabel}[1]{\quad#1.}
\makeatother



% Header and Footer with logo
\usepackage{lastpage,fancyhdr,graphicx}
\usepackage{epstopdf}
%\pagestyle{myheadings}
\pagestyle{fancy}
\fancyhf{}
%\setlength{\headheight}{27.023pt}
%\lhead{\includegraphics[width=2.0in]{PLOS-submission.eps}}
\rfoot{\thepage/\pageref{LastPage}}
\renewcommand{\headrulewidth}{0pt}
\renewcommand{\footrule}{\hrule height 2pt \vspace{2mm}}
\fancyheadoffset[L]{2.25in}
\fancyfootoffset[L]{2.25in}
\lfoot{\today}

%% Include all user-defined macros below

% Tables 1 and 4 are longer than one page; PLOS allows long tables to span pages, and a float cannot.
\usepackage{longtable}


%% END MACROS SECTION


\begin{document}
\vspace*{0.2in}

% Title must be 250 characters or less.
\begin{flushleft}
{\Large
\textbf{[TITLE --- author decision A1]} % Please use "sentence case" for title and headings (capitalize only the first word in a title (or heading), the first word in a subtitle (or subheading), and any proper nouns).
}
\newline
% Insert author names, affiliations and corresponding author email (do not include titles, positions, or degrees).
\\
{[AUTHOR NAME 1]}\textsuperscript{1*},
{[AUTHOR NAME 2 IF APPLICABLE]}\textsuperscript{1}
\\
\bigskip
\textbf{1} [AFFILIATION 1]
\\
\bigskip

% Use the asterisk to denote corresponding authorship and provide email address in note below.
* [CORRESPONDING AUTHOR EMAIL]

\end{flushleft}
% Please keep the abstract below 300 words
\section*{Abstract}
Drinking-water monitoring programmes must decide what to measure next, where and with which instrument, yet archives record where testing happened, not where risk lies; model evaluations can overstate performance in new places; and instrument documentation can be incomplete. We developed a procedure that passes this evidence through validation, permission and evidence gates and returns, for each monitoring context, an action and the smallest evidence that would change it. As a proof of principle, we applied it to 26 monitoring contexts in India, Kwale County (Kenya) and Malawi, and separately to published evidence from Bangladesh and Ghana. In India, further records could reclassify at least 17 of 160 eligible state or Union Territory \texttimes{} parameter cells at an assumed 0.10 action level for the share of records above a threshold. Models for seven tasks, fitted to a multinational, predominantly surface-water archive, averaged ROC-AUC 0.9396 under random-row splitting but 0.7696 under spatial block splitting. Separately fitted models, scored on Indian groundwater and a US well extract, kept their risk ordering for 5 of 13 archive \texttimes{} parameter pairs. Probabilities did not transfer: calibration slopes were 0.2383--0.4017 on those five (1 is ideal), Brier skill was negative on all six Indian pairs, calibration-in-the-large was not assessed and no local recalibration was run. Ordering could sequence confirmatory sampling; no transferred score could be read as a probability or replace a measurement. The 26 contexts received eight distinct actions, including repeat measurement and data repair. The procedure named no product, a single determination over 23 products: 891 of 1,242 product--criterion cells were not reported in any located document, and none documented all 11 required criteria. The Kenyan and Malawian archives are local, not national, and no model was used in Kenya, Malawi, Bangladesh or Ghana.

\clearpage
\newgeometry{top=0.85in,left=1in,right=1in,footskip=0.75in}
\linenumbers

\section*{Introduction}

Drinking-water quality monitoring reaches fewer people than it is meant to protect. Over the Sustainable Development Goal period, the share of the global population covered by total estimates of safely managed drinking water rose from 35\% to 72\%, and such estimates became available for India for the first time only in 2025~\cite{JMP2025}. Surveillance is also less frequent than required: in urban areas, 21 of 102 countries reported drinking-water surveillance at 95--100\% of the required frequency~\cite{GLAAS2025}. Infrequent sampling can overestimate how safe water is~\cite{Bain2014}. A programme with a fixed testing budget therefore has to decide which parameters to measure, at which water points and with which instrument, on evidence that is thin and unevenly spread.

Public archives record where and when testing happened, not the risk across a population. In India, the Central Ground Water Board's current monitoring framework samples every quality station once every five years before the monsoon and at least a quarter of the stations before and after the monsoon every year~\cite{CGWB-SOP-GWQDataAnalysis}. An exceedance rate computed from records collected under any such design is a property of the sampling as much as of the water. Measuring where the answer is already settled spends the budget without changing an action, while not measuring where the evidence is ambiguous leaves an action resting on ambiguous evidence.

One response is to model where contamination is likely. Predictive modelling can relate the probability that a contaminant such as groundwater arsenic exceeds a guideline value to soil, climate and terrain predictors, producing maps meant to guide further testing rather than replace it~\cite{Podgorski2020}. National models for India have mapped groundwater fluoride from 12,600 measurements~\cite{Podgorski2018} and groundwater arsenic, with the hazard maps proposed to inform the prioritisation of testing~\cite{Podgorski2020-IJERPH-India}.

How such models are evaluated decides what they can be used for, and practice differs. Environmental records are spatially autocorrelated, so a split that separates rows at random estimates something other than performance in an unseen region~\cite{Meyer2019}, and revisited sites compound the problem. Ignoring dependence between nearby samples biases model assessment towards optimism~\cite{MeyerPebesma2022}; the national Indian models were evaluated on random training and test splits~\cite{Podgorski2018,Podgorski2020-IJERPH-India}. Discrimination and calibration are likewise different properties: a model can order locations well while its stated probabilities are distorted~\cite{NiculescuMizil2005}. Estimated risks can be unreliable even when discrimination is good~\cite{VanCalster2019}, so a model whose ordering holds in a new setting may still misstate probabilities there.

Technology evidence raises the same question at the next step. Technology assessment characterises what instruments can measure and what deploying them demands, from power supply in remote settings to cleaning, signal drift and cost~\cite{Andres2018}. A programme needs performance near the threshold at which it must act. A 2003 verification of one arsenic field kit, in which one technical and one non-technical operator analysed identical samples, found false-negative rates of 62\% and 33\% respectively when they read its colour charts, relative to 10 \textmu{}g/L~\cite{EPA-ETV-QuickLowRangeII-2003}. A microbial water-quality test has been costed at 21.0 \textpm{} 11.3 USD on average across 18 African monitoring institutions in six countries, once consumables, equipment, labour and logistics are counted, with the \textpm{} not defined in the source~\cite{Delaire2017}.

Decision analysis makes trade-offs between objectives explicit and can carry uncertainty through to the comparison of alternatives~\cite{Schuwirth2018}; probabilistic forms report how probable each rank order is~\cite{Broekhuizen2015}. A ranking method returns a leader whether or not the criteria behind it are documented. Machine-learning models always make a prediction unless they are designed to abstain~\cite{Hendrickx2024}.

This article addresses the hand-offs between these bodies of evidence: what a programme should do when the evidence at one step is not strong enough to license the next. The hand-off is a recognised problem: a review of machine-learning prediction of United States drinking-water contaminants named accounting for spatial autocorrelation among the remaining methodological challenges~\cite{Hu2023}. The nearest existing approaches are the fixed surveillance schedules and the model-guided prioritisation of testing cited above; the procedure described here adds an explicit test of whether the evidence at each step is strong enough for the next, and a named output when it is not. Our aim was to assess, as separate steps of one decision, whether validation design, rank transfer, calibration under transfer and product documentation each support the use a monitoring programme would make of them, and so to state what to measure next, with what, and when the evidence does not support a decision. We applied the procedure as a proof-of-principle demonstration on archived and published evidence: to India as the principal application, to bounded local archives from Kwale County in Kenya and from Malawi, and to decisions in Bangladesh and Ghana from published evidence alone. It would be shown to be underspecified if two analysts applying its rules to the same archive returned different actions.

The aim was met in part. An action for every context holds by construction, because the rule returns one for any data state, so what the application tests is the evidence the gates consume. The procedure returned an action and the smallest decision-changing evidence for each of 26 monitoring contexts in India, Kenya and Malawi, and named no product in any of them, a single result from the 23-product documentary catalogue that applies to all 26. It did not license any use of a model probability, because calibration was poor on every transfer pair and local recalibration was not run. It was not compared with an alternative procedure.

\section*{Materials and methods}

\subsection*{The decision procedure}

The procedure takes three bodies of evidence (an accepted monitoring archive, the output of any predictive model, and the documentary evidence about candidate instruments) and returns two things for each monitoring context (Fig~\ref{fig1}A). The first is an action, which names what a programme should do next rather than scoring a place. Actions run from acquiring a baseline or repairing unusable records, through prioritising confirmatory sampling, to direct or targeted reference measurement and sentinel monitoring paired with chemistry; Table~\ref{table1} lists the action returned for each context and \nameref{S2_Table} defines all eight. The second return is the smallest decision-changing evidence, the least additional evidence that could move the context to a different action. Three gates stand between the inputs and the output, and a context meets them in order: the validation gate at step 6 of Fig~\ref{fig1}A, the permission gate at step 7 and the evidence gate at step 8.

\begin{figure}[!h]
\caption{\textbf{The procedure and the actions it returned.} (A) The workflow in four phases and eight steps. Evidence strands and populations are kept apart (step 1) and monitoring records are harmonised (step 2); accepted records are quantified (step 3) while documentary evidence for 23 exact products is classified against 54 criteria (step 4). Model output meets the validation gate, under three validation designs and on 13 external archive \texttimes{} parameter pairs (steps 5 and 6), and the permission gate, which keeps rank ordering apart from probability (step 7). The evidence gate requires all 11 active criteria to be documented before a product is named (step 8). The permitted outputs are the kinds of action the procedure may return, not the eight actions of panel B and not results. MCDA, multi-criteria decision analysis; CGWB, Central Ground Water Board; GEMStat, the global water-quality database of the UNEP GEMS/Water Programme. (B) The actions returned in the 26 monitoring contexts (India 8, Kenya 8, Malawi 10), each a country \texttimes{} source archive \texttimes{} parameter combination. No context named a product: that is one determination over the 23-product catalogue, applied to every context, not 26 determinations. The Kenyan and Malawian contexts come from bounded local archives and support no national statement. Plotted values: \nameref{S9_Data}.}
\label{fig1}
\end{figure}

\nolinenumbers
{\footnotesize\setlength{\tabcolsep}{3pt}\setlength{\LTcapwidth}{\linewidth}\captionsetup{font=normalsize}
\begin{longtable}{>{\raggedright\arraybackslash}p{1.55in}>{\raggedleft\arraybackslash}p{0.5in}>{\raggedright\arraybackslash}p{1.35in}>{\raggedright\arraybackslash}p{2.77in}}
\caption{\textbf{The action and the smallest decision-changing evidence returned for each of the 26 monitoring contexts.} One row per country \texttimes{} source archive \texttimes{} parameter combination. The smallest decision-changing evidence is the least additional evidence a programme would need for the procedure to return a different action; sample numbers, frequency, reference method and cost are not set by the procedure. No context names a product: that is one determination over the 23-product catalogue, applied to every context (note 1). The full matrix is \nameref{S2_Table}.}\label{table1} \\
\hline
\textbf{Monitoring context} & \textbf{Records (n)} & \textbf{Action} & \textbf{Smallest decision-changing evidence} \\ \thickhline
\endfirsthead
\hline
\textbf{Monitoring context} & \textbf{Records (n)} & \textbf{Action} & \textbf{Smallest decision-changing evidence} \\ \thickhline
\endhead
India, CGWB Year Book: fluoride & 35,843 & Rank prioritisation plus direct confirmation & Repair longitudinal well identity and run same-sample field-versus-reference validation. \\ \hline
India, CGWB Year Book: nitrate & 34,642 & Rank prioritisation plus direct confirmation & Repair persistent well/date/sample identity and collect repeat measurements before any temporal model. \\ \hline
India, CGWB portal: arsenic & 9,927 & Targeted direct reference measurement & Recover method and detection-limit metadata and build a defensible detected/non-detected analytical population. \\ \hline
India, CGWB portal: conductivity & 173,029 & Rank prioritisation plus direct confirmation & Governance review of the 1,500 \textmu{}S/cm project convention plus prospective sentinel-versus-spot evaluation; major-ion chemistry at flagged wells. \\ \hline
India, CGWB portal: pH & 162,729 & Direct measurement required & Collect representative local labels at sentinel sites and assess whether a separate local recalibration or model is useful; pH departures reported by direction. \\ \hline
India, CGWB portal: TDS & 30,151 & Rank prioritisation plus direct confirmation & Prospective paired conductivity/TDS sampling across relevant water matrices and seasons; recovery of how the portal TDS values were derived. \\ \hline
India, no accepted archive: \textit{E.~coli} & -- & Monitoring baseline required & Acquire traceable India microbial measurements with household/source context and compatible denominators. \\ \hline
India, no accepted archive: turbidity & -- & Monitoring baseline required & Acquire a machine-readable India archive with site, date, method, unit and valid turbidity results. \\ \hline
Kenya, Kwale groundwater samples: arsenic & 140 & Local signal plus direct confirmation & Repeat same-site arsenic measurement across additional seasons and confirm threshold classifications with a traceable reference workflow. \\ \hline
Kenya, Kwale groundwater samples: conductivity & 140 & Local signal plus direct confirmation & Repeat same-site conductivity measurement across additional seasons and confirm threshold classifications with a traceable reference workflow. \\ \hline
Kenya, Kwale groundwater samples: \textit{E.~coli} & -- & Monitoring baseline required & Acquire traceable local \textit{E.~coli} records with site, date, unit, method and quality-control fields. \\ \hline
Kenya, Kwale groundwater samples: fluoride & 139 & Local signal plus direct confirmation & Repeat same-site fluoride measurement across additional seasons and confirm threshold classifications with a traceable reference workflow. \\ \hline
Kenya, Kwale groundwater samples: nitrate & 139 & Local signal plus direct confirmation & Repeat same-site nitrate measurement across additional seasons and confirm threshold classifications with a traceable reference workflow. \\ \hline
Kenya, Kwale groundwater samples: pH & 140 & Local signal plus direct confirmation & Repeat same-site pH measurement across additional seasons and confirm threshold classifications with a traceable reference workflow. \\ \hline
Kenya, Kwale groundwater samples: TDS & -- & Monitoring baseline required & Acquire traceable local TDS records with site, date, unit, method and quality-control fields. \\ \hline
Kenya, Kwale groundwater samples: turbidity & -- & Monitoring baseline required & Acquire traceable local turbidity records with site, date, unit, method and quality-control fields. \\ \hline
Malawi, laboratory snapshot: arsenic & 10 & Data repair required & Recover arsenic units, method, LOD/LOQ and non-detect semantics before threshold analysis. \\ \hline
Malawi, laboratory snapshot: conductivity & 31 & Data repair required & Recover the conductivity scale/unit convention and raw instrument export; do not repair by assumption. \\ \hline
Malawi, laboratory snapshot: \textit{E.~coli} & -- & Organism-specific baseline required & Acquire organism-specific \textit{E.~coli} measurements with sample volume, reference method and response pathway. \\ \hline
Malawi, laboratory snapshot: faecal coliform & 19 & Organism-specific baseline required & Acquire organism-specific \textit{E.~coli} measurements with sample volume, reference method and response pathway. \\ \hline
Malawi, laboratory snapshot: fluoride & 2 & Data repair required & Recover the fluoride missing-value codebook and valid observations before threshold analysis. \\ \hline
Malawi, laboratory snapshot: nitrate & 31 & Local signal plus direct confirmation & Expand repeat-site nitrate coverage and retain method, calibration, operator and quality-control metadata. \\ \hline
Malawi, laboratory snapshot: pH & 31 & Local signal plus direct confirmation & Expand repeat-site pH coverage and retain method, calibration, operator and quality-control metadata. \\ \hline
Malawi, laboratory snapshot: TDS & 31 & Local signal plus direct confirmation & Expand repeat-site TDS coverage and retain method, calibration, operator and quality-control metadata. \\ \hline
Malawi, laboratory snapshot: turbidity & 22 & Local signal plus direct confirmation & Expand repeat-site turbidity coverage and retain method, calibration, operator and quality-control metadata. \\ \hline
Malawi, logger series: conductivity & 714 & High-frequency sentinel plus chemical confirmation & Pair the two persistently high-conductivity sites and comparison sites with direct chemistry, calibration, fouling, downtime and energy records. \\ \hline
\end{longtable}
\begin{flushleft}
1. Every context returns an action and the smallest decision-changing evidence. \textbf{Every context also abstains at product level}: the procedure names no product in any of the 26, because no product's located documentation reports every criterion the decision needs. That abstention is one determination over the 23-product catalogue, applied to every context, and is independent of which action a context returns.

2. \textbf{An en dash means the record count is not available}: no accepted measurement population exists for that context in the archives analysed. It does not mean zero exceedances, and absence of records is not evidence of safety. The six are India \textit{E.~coli} (no accepted India measurement population); India turbidity (no accepted India measurement population); Kenya \textit{E.~coli} in the Kwale groundwater samples (data not available); Kenya TDS in the Kwale groundwater samples (data not available); Kenya turbidity in the Kwale groundwater samples (data not available); Malawi \textit{E.~coli} in the laboratory snapshot (data not available; faecal coliform is not \textit{E.~coli}). The procedure does not specify whether the evidence needed is existing records from another custodian or new measurement.

3. Kwale County and the two Malawi archives are bounded local archives. \textbf{None is a national estimate}, and none supports a country-level statement.

4. Faecal coliform is reported as its own context and is never relabelled as \textit{E.~coli}.

5. The records column counts accepted records, except for Malawi conductivity, fluoride and arsenic, which failed acceptance, for which it counts the rows held, and for the logger series, for which it counts site-days.

6. In the Indian conductivity, TDS and pH rows, the clause after the last semicolon of the evidence column was added for this table and is not in the context matrix (\nameref{S2_Table}), which keeps the recorded text.

7. The evidence column states what would change an action. The procedure does not set who supplies that evidence, the number of samples or sites, the frequency, the reference method or the cost, nor the action a context would move to. Data repair can be made only by the originating laboratory or data publisher, and the procedure has no fallback if records cannot be recovered. Each context's operational requirements, including incubation energy and cost per valid sample, are listed in \nameref{S2_Table} and were not costed.

8. In Kwale, two of the three fluoride exceedances and two of the 19 conductivity exceedances come from one thermal spring, and the third fluoride exceedance is an unverified value. In the Malawi snapshot, TDS is recorded on mixed scales in 10 of 31 records, as conductivity is, and the zero nitrate count is inconclusive rather than a clearance. The actions are shown as the rule returned them.

9. \textit{Local labels} are direct measurements used to check a model's scores; \textit{abstain} means that no product is named; LOD/LOQ are the limits of detection and quantification. For Indian conductivity, the \textit{governance review} means replacing the project convention with the action value a programme uses, none of which was assessed. The procedure does not name the reference method meant by a \textit{traceable reference workflow}.
\end{flushleft}
}
\linenumbers

The validation gate settles in which role a model may enter a decision: ordering locations so that confirmatory sampling is sequenced, supplying a probability to compare with an action level, or neither. The role follows from the validation design, not from the size of a reported metric. Monitoring records are not exchangeable rows, because sites are revisited and nearby observations share geology, land cover and often a sampling programme, so a random split mostly measures interpolation within the existing sampling structure. Keeping every record from a site together removes the repeat-visit component; holding out contiguous spatial blocks additionally requires the model to generalise across geography~\cite{Meyer2019}. The gate stops a random-row figure from being read as performance in a district that contributed no training station.

Discrimination asks whether the model orders cases correctly; calibration asks whether a stated probability matches the frequency actually observed at that probability~\cite{NiculescuMizil2005}. Calibration here always means model calibration, not the calibration of an instrument. A ranking tells a programme where to look first, and a probability compared with a threshold tells it when to act. The permission gate therefore recognises three states. Where ordering survives a move to a new setting but calibration there has not been shown to be adequate, the model may set the order of confirmatory measurement and nothing more. Where ordering survives and calibration on local labels from that setting has been assessed and found adequate, or restored by local recalibration, a probability may also be compared with a fixed action level. Where ordering does not survive, the model contributes nothing and direct measurement is required. No adequacy criterion for calibration was fixed in advance in this study. The gate is applied to each archive \texttimes{} parameter pair, because two archives holding the same analyte can return opposite verdicts.

Once a decision requires direct measurement it becomes a question about an instrument, and the evidence gate governs it through two rules. The documentation threshold asks whether a criterion is documented for enough products for a comparison on it to mean anything. The sufficiency rule asks whether a given product documents every active criterion and passes its hard analytical, matrix, deployment and identity requirements. A product that fails the sufficiency rule is not scored, and the procedure returns a product-level abstention together with the documents that would lift it. Scoring an undocumented field at a middling value would reward weak disclosure, because a device that reports a poor maintenance interval would then rank below one that reports nothing. Missing documentation is unknown evidence about a product, not evidence that it performs poorly. Multi-criteria comparison of technology archetypes gives preference context for a class of instrument under stated weights, is not screened for the analytes a context needs, and establishes nothing about any product within that class.

We distinguish three monitoring functions, because comparing devices before the function is fixed ranks alternatives that were never doing the same job. Distributed spot screening covers existing water points. A higher-frequency sentinel at selected sites buys temporal resolution at the cost of power, logging, drift, fouling and downtime. Confirmatory or reference measurement establishes a contaminant-specific result. The functions combine within one monitoring architecture, no technology class is presumed to serve all three, and each setting's action names the function it relies on.

Abstention is accordingly a designed output with the same role as the reject option of a selective classifier~\cite{Geifman2017}, but triggered by missing evidence rather than by expected error (see Discussion), and not a failure. Because the gates are ordered, a context whose records cannot be accepted returns a baseline or repair action before any model is considered, one whose model output fails the validation or permission gate returns direct measurement, and one that clears both still reaches the evidence gate before any instrument is named. The rules are the same whatever a product documents, so no product's values can alter the gates it is judged by. How the gates were implemented for the contexts reported here, including which outcomes the authors recorded rather than computed, is set out below, under Technology evidence and actions.

\subsection*{Design, settings and unit of analysis}

The unit of analysis for decisions is the monitoring context, one country \texttimes{} source archive \texttimes{} parameter combination, and 26 were evaluated, 8 in India, 8 in Kenya and 10 in Malawi (\nameref{S1_Text} gives how they were generated). The Kenyan and Malawian archives are bounded local archives, not national samples, and no result from them is a national estimate for either country. Bangladesh and Ghana are applications of the decision logic to published evidence, with no model fitted or transferred, and their studies were chosen for the decision each setting poses rather than found by a systematic search. Five populations are kept separate and never merged: 23 exact commercial products, a 33-entry technology landscape (Fig~\ref{fig1}A), reference laboratory methods, emerging research classes, and ten technology archetypes, classes of instrument compared by multi-criteria decision analysis (MCDA). Missing evidence was retained, not imputed.

\subsection*{Archives}

Six evidence strands enter the analysis (Fig~\ref{fig2}): the fitting archive, the Indian archives, the United States archive, the Kenyan and Malawian archives, published studies from Bangladesh and Ghana, and documentation for named products. The fitting archive is the GEMStat global water-quality database of the UNEP GEMS/Water Programme~\cite{GEMStat}: a multinational archive, predominantly of surface water, used only for model fitting and internal evaluation. Seven classification tasks were built from it, for arsenic, nitrate, conductivity, suspended solids, turbidity, fluoride and pH (thresholds in \nameref{S1_Text}).

\begin{figure}[!h]
\caption{\textbf{The evidence strands and where they come from.} (A) The role of each country's evidence. India holds the two Central Ground Water Board (CGWB) external archives, portal records and Year Book tables that share a producer, and is the principal application; the United States Water Quality Portal well-site extract, from four states, is the second external arm; Kenya and Malawi are bounded local archives; and Bangladesh and Ghana enter as applications of the decision logic to published evidence. Shading shows a country's role, not the extent of sampling. Station locations of GEMStat, the multinational and predominantly surface-water fitting archive, and of the CGWB archives are not drawn, because their redistribution terms have not been established. (B) The 81 sampling locations of the Kwale County archive, from two 2016 campaigns~\cite{OWD-grdwtrsmpkwale}, reprojected from UTM zone 37S on the Arc 1960 datum without a datum shift; no coastline is drawn at this scale. Inset: Kenya. (C) The laboratory snapshot of 32 Malawian water points~\cite{OWD-boreholelabdata} and the conductivity logger series at 19 sites~\cite{OWD-mwgroundwaterdata}. Nearby markers overlap at this scale. Country boundaries: Natural Earth (public domain). The openwashdata packages are published under CC BY 4.0. Plotted values: \nameref{S9_Data}.}
\label{fig2}
\end{figure}

Two external arms were used to evaluate transfer, and neither took part in fitting. The primary arm is Indian groundwater from the Central Ground Water Board (CGWB), in two products that share a producer: National Water Data Portal records for pH, conductivity, total dissolved solids (TDS) and arsenic (the CGWB portal), and fluoride and nitrate tables recovered from the CGWB Ground Water Quality Data documents (the CGWB Year Book)~\cite{CGWB-NWDP,CGWB-GWQ2019,CGWB-GWQ2021,CGWB-GWQ2022}. CGWB samples groundwater from observation and monitoring wells, with the stated aim of assessing shallow aquifers against the Indian drinking-water standard~\cite{CGWB-AGWQR2025}, so a record describes groundwater at a monitoring well, not water as supplied. The second arm, WQP, is a 2023 well-site extract for four states from the United States Water Quality Portal~\cite{WQP} (\nameref{S1_Text}). Both external arms are groundwater, so transfer from the fitting archive crosses water-body type as well as geography. Independence rests on archive lineage and a 0.1\textdegree{} cell exclusion between training and external locations, not on a record-level audit, so no fully external claim is made.

The bounded local archives are three openwashdata packages published under CC BY: samples from two 2016 campaigns in Kwale County, Kenya~\cite{OWD-grdwtrsmpkwale}; a laboratory snapshot of 32 Malawian water points, 30 of them boreholes or tubewells, sampled in three rounds with different parameter suites~\cite{OWD-boreholelabdata}; and a Malawian conductivity series from loggers in monitoring wells~\cite{OWD-mwgroundwaterdata}. Surface-water samples were excluded from the Kwale archive, which records no water-point type. Acceptance checked units, identity and plausibility, and harmonisation rules, versions and accepted record counts are in \nameref{S1_Text}.

\subsection*{Thresholds and chemical basis}

An exceedance rate is a record fraction, the proportion of accepted records above a threshold, reported with its numerator and denominator (Table~\ref{table2}). Fluoride above 1.5 mg/L and arsenic above 0.010 mg/L follow the WHO guideline values, the arsenic value being provisional~\cite{WHO2022GDWQ}. These WHO values are applied to every archive, whereas CGWB judges Indian drinking suitability against the Indian standard IS 10500~\cite{CGWB-AGWQR2025}, whose limits \nameref{S1_Text} compares with ours. Arsenic values labelled mg/L but implausible as such, from four north-eastern states in 2017, were divided by 1,000 (\nameref{S1_Text}). Nitrate was classified above 50 mg/L as NO\textsubscript{3}, the WHO guideline value expressed as nitrate ion~\cite{WHO2022GDWQ}. The transfer tasks use nitrate above 10 mg/L as N, the value of the United States maximum contaminant level~\cite{USEPA-NPDWR}. That is 44.268 mg/L as NO\textsubscript{3}, so no nitrate figure is carried between the two analyses. WHO proposes no health-based guideline value for pH or TDS: 6.5--8.5 is the usual optimum pH range for supply, and water becomes increasingly unpalatable at TDS above about 1,000 mg/L~\cite{WHO2022GDWQ}; both are used here as operational ranges. Conductivity above 1,500 \textmu{}S/cm is an unverified operational project convention, not a health guideline or an Indian drinking-water limit; under the factor CGWB recommends, TDS of about 0.65 \texttimes{} conductivity~\cite{CGWB-SOP-GWQDataAnalysis}, it corresponds to about 975 mg/L TDS. The 1 NTU turbidity rule is WHO's value for supporting effective disinfection~\cite{WHO2022GDWQ}, applied here to untreated water as a project rule. Microbial results use a 100 mL basis, and faecal coliform is never relabelled as \textit{E.~coli}.

\begin{table}[!ht]
\centering
\caption{\textbf{The Indian evidence base by parameter, including the cells whose Wilson intervals straddle a 0.10 action level.} One row per parameter, aggregated from 186 state or Union Territory \texttimes{} parameter cells. Exceedance rates are record fractions with Wilson 95\% intervals that treat records as independent. The notes give the source archive of each parameter, the states with data, and the median revisit interval, which is recorded in 145 cells and is annual or longer in all of them.}
{\footnotesize\setlength{\tabcolsep}{3pt}
\begin{tabular}{>{\raggedright\arraybackslash}p{0.72in}>{\raggedleft\arraybackslash}p{0.52in}>{\raggedleft\arraybackslash}p{0.55in}>{\raggedleft\arraybackslash}p{0.48in}>{\raggedright\arraybackslash}p{0.80in}>{\raggedright\arraybackslash}p{1.15in}>{\raggedleft\arraybackslash}p{0.55in}>{\raggedleft\arraybackslash}p{0.50in}>{\raggedleft\arraybackslash}p{0.48in}}
\hline
\textbf{Parameter} & \textbf{Accepted records (n)} & \textbf{Records above threshold (n)} & \textbf{Unique sites (n)} & \textbf{Exceedance threshold} & \textbf{Record fraction [95\% Wilson interval]} & \textbf{Cells supporting description (of 31)} & \textbf{Decision-relevant cells (n)} & \textbf{Cells supporting leave-one-state-out (of 31)} \\ \thickhline
Conductivity & 173,029 & 39,115 & 24,128 & \textgreater{}1,500 \textmu{}S/cm (operational convention) & 0.2261 [0.2241, 0.2280] & 30 & 1 & 22 \\ \hline
pH & 162,729 & 10,409 & 23,704 & outside 6.5--8.5 pH units & 0.0640 [0.0628, 0.0652] & 30 & 4 & 21 \\ \hline
Fluoride & 35,843 & 2,344 & 32,561 & \textgreater{}1.5 mg/L & 0.0654 [0.0629, 0.0680] & 27 & 6 & 14 \\ \hline
Nitrate & 34,642 & 6,403 & 32,206 & \textgreater{}50 mg/L as NO\textsubscript{3} & 0.1848 [0.1808, 0.1890] & 27 & 2 & 18 \\ \hline
TDS & 30,151 & 5,442 & 14,632 & \textgreater{}1,000 mg/L & 0.1805 [0.1762, 0.1849] & 27 & 4 & 14 \\ \hline
Arsenic & 9,927 & 31 & 7,931 & \textgreater{}0.010 mg/L & 0.0031 [0.0022, 0.0044] & 19 & 0 & 0 \\ \hline
\end{tabular}
\begin{flushleft}
1. \textbf{Median revisit interval.} In every cell with a recorded median, the median interval between successive records at a site is 365 days or longer: all 145 of 145 (conductivity 31; pH 31; fluoride 23; nitrate 23; TDS 30; arsenic 7), and \textbf{48 are exactly 365 days}. Of these cells 99 come from the portal and 46 from the Year Book tables, whose medians fall only at edition spacings (365, 731, 1,096 days) and cannot fall below them. 17 medians rest on fewer than five intervals, the fewest on 1. The interval describes the archive as held and prescribes no sampling frequency.

2. A cell is \textbf{decision-relevant} when its Wilson interval on the exceedance rate straddles an assumed 0.10 action level, so that more records could establish on which side it lies: \textbf{17 of 160} cells that support description. The intervals treat records as independent although sites are revisited, so this count is a lower bound. The classification marks where evidence could change a state-level classification. It is not severity, prevalence, exposure or a hotspot ranking; a cell resolved below the level is not one where monitoring may stop; and a state or Union Territory cell is not a well or a person.

3. The last column counts the cells whose record support would permit a leave-one-state-out evaluation, the design that would test whether a model trained on some states carries to a state it has never seen. \textbf{For arsenic that count is zero.} Support for description is therefore not support for generalisation, and the two columns must never be read as one.

4. The \textbf{1,500 \textmu{}S/cm} conductivity value is an operational project convention, not a health guideline, and no health inference follows from exceeding it.

5. The nitrate threshold is 50 mg/L \textbf{as NO\textsubscript{3}}, not as N. The transfer analysis uses 10 mg/L as N, the United States maximum contaminant level, a different rule, so nitrate exceedance figures are not carried between that analysis and this table.

6. Fluoride and nitrate come from the Year Book tables; the other four from the National Water Data Portal. \textbf{The two routes carry different rights positions}, so no single licence statement covers `the Indian data'.

7. Intervals are Wilson score intervals on the exceedance proportion. The arsenic action in Table~\ref{table1} follows the evidence state recorded for arsenic (insufficient transfer support and a denominator-sensitive rate); separately, a unit correction was applied upstream to a subset of arsenic records (\nameref{S1_Text}).

8. Records come from all 31 states and Union Territories for conductivity, pH, fluoride, nitrate and TDS; arsenic records come from 23 of them. Per-cell counts, intervals and revisit intervals are in \nameref{S3_Table}.
\end{flushleft}
}
\label{table2}
\end{table}

Binomial uncertainty on record-level exceedance proportions is described by Wilson 95\% intervals~\cite{Brown2001}, which do not remove spatial clustering or sampling-design bias, so rates are reported as monitoring-record signals, not as population prevalence or exposure.

\subsection*{Analyses of the settings}

For India we formed state or Union Territory \texttimes{} parameter cells from the accepted CGWB records. A cell was described when it held at least 30 records at 10 or more sites; its median revisit interval, on a site key built from the station or well identifier and the sample date, and its support for leave-one-state-out and temporal analysis were recorded (\nameref{S3_Table}). A cell is decision-relevant when the Wilson 95\% interval on its exceedance rate straddles an action level, set at 0.10 as a stated assumption and not a regulatory value, with counts also reported at 0.05 and 0.25. This classification allocates confirmatory effort between states and parameters and sets no action in Table~\ref{table1}; CGWB's own response to an exceedance is at station level, monitoring affected areas as trend wells and new hotspots on a 2 km grid~\cite{CGWB-SOP-GWQDataAnalysis}. For Kwale and the Malawi laboratory snapshot we report record fractions against the same thresholds. For the 67 Kwale sites sampled in both campaigns (66 for nitrate and fluoride) we compared threshold states between campaigns using site medians, and for the Malawi logger series we counted changes of state between successive daily site medians.

\subsection*{Models and validation}

Model A uses environmental and terrain covariates and the site coordinate, and Model B uses earlier measurements at the same site and needs at least two of them. Both use the histogram-based gradient-boosting classifier in scikit-learn~\cite{Pedregosa2011}, with configurations fixed in advance and no hyperparameter search. Model A's configuration is in \nameref{S1_Text}; Model B's was not retained.

Three validation designs were compared on the seven fitting tasks (Fig~\ref{fig3}, Table~\ref{table3}). Each split the fitting archive once into training, validation and test partitions keyed on the record (random-row), the station (station-grouped) or a contiguous 3\textdegree{} spatial block; Model A was fitted on 24 environmental, terrain and coordinate features, calibrated isotonically on the validation partition and scored once on the test partition (\nameref{S1_Text}). Block splitting keeps nearby observations together and is the design used here to address generalisation to held-out blocks within the sampled geography~\cite{Meyer2019}; no autocorrelation range was estimated to set the block size. ROC-AUC (area under the receiver operating characteristic curve), precision--recall AUC, sensitivity, specificity, calibration slope and Brier skill were interpreted separately; precision--recall AUC is read against prevalence as its no-skill baseline~\cite{Saito2015}, and a 0.5 cut-point was used diagnostically rather than as an operational default. Prevalence here is a task's base rate, the share of its records above threshold, not population prevalence. A calibration slope between 0 and 1 means predicted probabilities are too extreme, and a slope at or below 0 means they carry no usable probability ordering. The two channels were compared separately, in forward-chained yearly folds on records with at least two earlier same-site measurements, with DeLong tests on the pooled predictions (\nameref{S1_Text}, Table~\ref{table3}). No local recalibration was performed, because it needs held-out predictions that were not retained.

\begin{figure}[!h]
\caption{\textbf{Discrimination of the fitting-archive models under three validation designs, and the paired channel comparison.} (A) Schematic of the designs; positions carry no geographic meaning. Random-row splitting can place records from one station on both sides of the split, station-grouped splitting keeps each station whole, and spatial-block splitting assigns contiguous 3\textdegree{} blocks. (B) Test-partition ROC-AUC for the seven classification tasks, with means of 0.9396 (random-row), 0.8688 (station-grouped) and 0.7696 (spatial-block), an optimism of 0.1701. The random-row value measures leakage by difference, and no design supports transfer to an unmonitored country. (C) Spatial-block PR-AUC against each task's prevalence, its no-skill baseline. (D) ROC-AUC of Model A (environmental, positional and neighbour features) and Model B (antecedent measurements at the same site) for the six tasks with a paired comparison, in forward-chained yearly folds on records with at least two earlier same-site measurements; this comparison is not evidence of spatial generalisation. Values are from retained result tables (Table~\ref{table3}). ROC-AUC, area under the receiver operating characteristic curve; PR-AUC, area under the precision--recall curve. Plotted values: \nameref{S9_Data}.}
\label{fig3}
\end{figure}

\begin{table}[!ht]
\centering
\caption{\textbf{Discrimination and calibration of the fitting-archive models under three validation designs, and the paired channel comparison.} One row per parameter. Prevalence and test-partition ROC-AUC under random-row, station-grouped and spatial-block splitting; optimism is random-row minus spatial-block ROC-AUC, and PR-AUC and the calibration slope are for the spatial-block design. The last three columns compare Model A and Model B in forward-chained yearly folds on records with at least two earlier same-site measurements, which is not evidence of spatial generalisation.}
{\footnotesize\setlength{\tabcolsep}{3pt}
\begin{tabular}{>{\raggedright\arraybackslash}p{0.66in}>{\raggedleft\arraybackslash}p{0.62in}>{\raggedleft\arraybackslash}p{0.44in}>{\raggedleft\arraybackslash}p{0.49in}>{\raggedleft\arraybackslash}p{0.44in}>{\raggedleft\arraybackslash}p{0.60in}>{\raggedleft\arraybackslash}p{0.44in}>{\raggedleft\arraybackslash}p{0.60in}>{\raggedleft\arraybackslash}p{0.43in}>{\raggedleft\arraybackslash}p{0.43in}>{\raggedleft\arraybackslash}p{0.43in}}
\hline
\textbf{Parameter} & \textbf{Prevalence} & \textbf{ROC-AUC random row} & \textbf{ROC-AUC station grouped} & \textbf{ROC-AUC spatial block} & \textbf{Optimism} & \textbf{Spatial PR-AUC} & \textbf{Spatial calibration slope} & \textbf{Model A ROC-AUC} & \textbf{Model B ROC-AUC} & \textbf{Model B \textminus{} A} \\ \thickhline
Arsenic & 0.0833 & 0.9626 & 0.9230 & 0.8160 & 0.1466 & 0.2710 & 0.8922 & 0.9302 & 0.9530 & +0.0228 \\ \hline
Nitrate & 0.0263 & 0.9727 & 0.8527 & 0.6800 & 0.2927 & 0.0743 & 0.8363 & 0.9444 & 0.9779 & +0.0335 \\ \hline
Conductivity & 0.0416 & 0.9854 & 0.9411 & 0.8596 & 0.1258 & 0.3428 & 0.7001 & 0.9633 & 0.9858 & +0.0225 \\ \hline
Suspended solids & 0.2277 & 0.8683 & 0.8203 & 0.7463 & 0.1221 & 0.4633 & 0.8541 & 0.8508 & 0.8724 & +0.0215 \\ \hline
Turbidity & 0.8323 & 0.9293 & 0.8615 & 0.7440 & 0.1853 & 0.9272 & 0.6749 & 0.9121 & 0.9426 & +0.0305 \\ \hline
Fluoride & 0.0228 & 0.9830 & 0.9336 & 0.9062 & 0.0768 & 0.1912 & 0.1925 & -- & -- & -- \\ \hline
pH & 0.0715 & 0.8760 & 0.7496 & 0.6348 & 0.2412 & 0.1855 & 0.2979 & 0.8298 & 0.8940 & +0.0643 \\ \hline
\textbf{Mean of 7} & n/a & \textbf{0.9396} & \textbf{0.8688} & \textbf{0.7696} & \textbf{0.1701} & n/a & n/a & n/a & n/a & n/a \\ \hline
\end{tabular}
\begin{flushleft}
1. \textbf{Optimism is random-row ROC-AUC minus spatial-block ROC-AUC.} The three designs answer three different questions and none of them is `the' performance of the model. Random-row splitting places records from one station on both sides of the split, and its values are reported \textbf{to measure leakage by difference}, not as an estimate of performance at new locations; spatial-block splitting is the only design that speaks to a new area. \textbf{None of the three is a claim about transfer to an unmonitored country.} Optimism and the Model B \textminus{} A differences are computed from unrounded values, so they can differ by 0.0001 from a difference of the rounded values shown.

2. Model A (environmental channel: the 24 validation features, eight neighbour features and month, 33 in all; a separate fit from the one behind the validation-design columns, \nameref{S1_Text}) and Model B (site-history channel: 12 antecedent-history features) were compared in \textbf{forward-chained yearly folds}, each year predicted from all earlier years on records with at least two earlier measurements at the site, so the comparison is not evidence of spatial generalisation, and the difference under station-grouped or spatial designs was not estimated. \textbf{Model B exceeds Model A on all six parameters that have a paired comparison}, both in discrimination and in closeness of the calibration slope to 1 (Model A against Model B: arsenic 0.7608 against 0.8883; nitrate 0.8284 against 0.9344; conductivity 0.7922 against 0.9603; suspended solids 0.8700 against 0.9695; turbidity 0.7689 against 0.9156; pH 0.8442 against 0.9570). The advantage requires prior measurements at the site, so it is unavailable exactly where monitoring is absent.

3. \textbf{DeLong p-values are not reported.} The frozen table records 0.00e+00 for all six paired comparisons, an exported value that the project's verification record describes as ``below display precision, not a literal exact zero'', and the tests run on 174,533 to 803,735 paired observations. The test assumes independent observations, which revisited stations violate, and at those sizes a significance test on a ROC-AUC difference is dominated by sample size and says nothing about practical importance. \textbf{The effect size in the `Model B \textminus{} A' column, which spans +0.0215 to +0.0643, is the quantity to read.}

4. \textbf{An en dash means not available}: fluoride has no paired Model A/B comparison in the frozen results. \textbf{`n/a' means not applicable}: a mean is not reported for prevalence, which is a property of the data rather than of performance, nor for the paired-comparison columns, which cover six parameters and are described in note 2. Prevalence here is a task's base rate, the share of its records above threshold, not population prevalence.

5. A mean is deliberately \textbf{not} reported for the spatial calibration slope, whose values span 0.1925 to 0.8922, nor for spatial PR-AUC, which spans 0.0743 to 0.9272. Averaging either would conceal the spread that matters. \textbf{PR-AUC must be read against the prevalence in the same row}: turbidity's PR-AUC of 0.9272 accompanies a prevalence of 0.8323, so for that parameter the majority class is the positive one.

6. \textbf{Calibration is reported here only as the spatial-block slope within the fitting archive}, where it was below 1 for all seven tasks. Within the fitting archive, spatial-block discrimination supports ordering by relative risk; ordering in any other archive rests on the transfer verdicts, and calibration under transfer is reported with them. These results do not support a fixed probability threshold, and a modelled cell is not a well, tap, source, household or person.

7. These values are cited from frozen result tables. \textbf{The runs behind them were not re-run for this article and no prediction-level file was retained}; the run scripts are held by the author, and verifying a result table is not reproducing a model.
\end{flushleft}
}
\label{table3}
\end{table}

A mechanistic ablation tested whether aquifer, soil, land-cover and population variables substitute for raw coordinates for conductivity and pH; its result, whose intervals overlap, is in \nameref{S5_Text}, and \nameref{S6_Fig} shows one illustrative surface, read as an ordering and not as a probability. Its stack comprises the national aquifer classification~\cite{CGWB-ManualAquiferMapping}, SoilGrids 2.0 soil properties~\cite{Poggio2021}, ESA WorldCover 10 m 2021 v200 land cover~\cite{Zanaga2022} and WorldPop population density~\cite{WorldPop2020}.

\subsection*{External transfer}

Transfer was assessed per archive \texttimes{} parameter pair. For each parameter a separate Model A fit was trained on GEMStat records before a cut year, on 33 features that add nearest-station neighbour features and month to the 24 above (\nameref{S1_Text}). These transfer fits are not the seven validation fits and were not evaluated under the three designs. The frozen fits were scored on the external archives with no tuning, threshold selection or recalibration after the external data were seen, once every external site sharing a 0.1\textdegree{} cell with a training station had been discarded. Of 24 possible pairs, 13 were evaluated and 11 are reported as untested rather than negative.

The verdict comes from three interval tests rather than from a threshold on discrimination: whether the site-level rank correlation between predicted risk and observed severity has a 95\% site-bootstrap interval wholly above 0, whether the rank correlation within 1\textdegree{} blocks has a block-bootstrap interval wholly above 0, and whether lift in the top predicted decile has a site-bootstrap interval wholly above 1. With fewer than two tests evaluable a pair has \textit{insufficient support}, a gap in the evidence that is never merged with a negative result; if every evaluable test passes, \textit{rank transfer is supported}; if none passes, it is \textit{not supported}; otherwise it is \textit{partial or scale-dependent}. Only the outcome of each test was retained, not its statistic or interval. For every pair we also recorded threshold discrimination and the Spearman rank correlation between predicted risk and observed severity among detected records, each with a site-clustered bootstrap interval, the calibration slope, Brier skill against a reference that knows the external base rate, and confusion counts at a 0.5 cut-point (\nameref{S1_Text}). The retained calibration intercept was estimated jointly with the slope and so is not calibration-in-the-large, and no criterion for adequate calibration was fixed in advance. A supported verdict authorises prioritised confirmatory sampling only, at the level of observations within the archive.

\subsection*{Technology evidence and actions}

The exact-product catalogue comprises 23 named commercial products assessed against 54 criteria, a 1,242-cell grid in which each cell records its value, source, locus and evidence class (\nameref{S4_Table}). Not reported, not researched, conflicting and range-only evidence are kept distinct, and evidence about an adjacent product cannot verify a named model. A criterion meets the documentation threshold when it is documented for at least 12 of the 23 products, counted on its primary database field (strict) or on any database field mapped to it (any admissible). The sufficiency rule requires all 11 active criteria (\nameref{S1_Text}) to be documented; threshold-near classification error, invalid-test rate and testing frequency are not among the active criteria.

Conditional comparison of the ten technology archetypes used the same 11 criteria under seven declared weighting scenarios, with 40,000 Monte Carlo draws per scenario from a Dirichlet distribution centred on the scenario's weights (\nameref{S1_Text}, \nameref{S7_Table}). The technique for order of preference by similarity to ideal solution (TOPSIS) supplied the primary rank~\cite{Chakraborty2022}, and stochastic multicriteria acceptability analysis (SMAA) represented weight uncertainty as a distribution over the feasible weights~\cite{TervonenLahdelma2007}. The comparison is not screened for the analytes a scenario targets, and no practitioner weights were elicited or invented.

For each monitoring context a rule returns an action (\nameref{S2_Table}): for Kenya and Malawi from the data status alone, whatever the exceedance count, and for India from a per-parameter evidence state that the authors assigned from the transfer verdict, record support and identity limits (\nameref{S1_Text} gives both branches). The evidence states and the rule were recorded on 13 August 2026, after the transfer verdicts had been computed, and no dated protocol records when the action level, support requirement, verdict rule, documentation threshold or active criteria were fixed. The product-level decision is a single determination over the 23-product catalogue, applied to all 26 contexts, not 26 independent determinations.

\subsection*{Reproducibility}

Machine-readable result exports were retained for the India and WQP transfers and the mechanistic ablation, and the canonical result tables were verified against them by a script that does not refit; the output file of the validation-design run matches Table~\ref{table3} to four decimals. The run scripts are held by the author and were not re-run for this article, and no prediction-level file was retained. A result-table verification is not a full training reproduction, and no result in this paper is described as reproduced, replicated or independently re-derived. The scikit-learn version used for the model runs was not recorded (\nameref{S1_Text}).

\subsection*{Artificial-intelligence assistance}

We used large language models in preparing this manuscript: OpenAI ChatGPT (version not recorded) to set up the initial project file structure of prompts, rules and scripts, and Anthropic Claude Opus 5, through Claude Code, to write and run analysis, verification and figure-building scripts, open cited sources and record the passage supporting each citation, and draft and revise prose. No AI tool altered or imputed a measurement, and every count, table and quoted value was checked by script against its source file or recorded passage. \nameref{S10_Text} records each activity, its check and whether the author has reviewed it; at the time of writing, the author's review of the prose and of the citation records was outstanding.

\section*{Results}

We applied the procedure first to India, where the accepted archives are largest, then to the bounded local archives of Kwale County in Kenya and of Malawi, and finally to published evidence from Bangladesh and Ghana (Fig~\ref{fig2}). Table~\ref{table4} sets out, for each setting, the evidence held, what the model may and may not do, the action returned and the evidence that would change it, and \nameref{S12_Table} adds the monitoring decision each setting poses and where the procedure declines.

\nolinenumbers
{\footnotesize\setlength{\tabcolsep}{3pt}\setlength{\LTcapwidth}{\linewidth}\captionsetup{font=normalsize}
\begin{longtable}{>{\raggedright\arraybackslash}p{0.86in}>{\raggedright\arraybackslash}p{1.64in}>{\raggedright\arraybackslash}p{1.17in}>{\raggedright\arraybackslash}p{1.17in}>{\raggedright\arraybackslash}p{1.25in}}
\caption{\textbf{For each application setting, the evidence held, what a model may do, the action returned and the evidence that would change it.} One row per setting: India; Kwale County, Kenya; the Malawi laboratory snapshot; the Malawi logger series; and Bangladesh and Ghana, whose evidence is published and to which no model was transferred. Counts above a threshold are record fractions within each archive, not prevalence. Table~\ref{table1} lists the 26 monitoring contexts behind the India, Kenya and Malawi rows, and \nameref{S12_Table} adds the monitoring decision each setting poses and where the procedure declines.}\label{table4} \\
\hline
\textbf{Setting} & \textbf{Evidence held} & \textbf{What a model may do} & \textbf{Action returned} & \textbf{Evidence that would change it} \\ \thickhline
\endfirsthead
\hline
\textbf{Setting} & \textbf{Evidence held} & \textbf{What a model may do} & \textbf{Action returned} & \textbf{Evidence that would change it} \\ \thickhline
\endhead
India (CGWB portal and CGWB Year Book) & 186 cells in 31 states and Union Territories; 160 eligible; 17 decision-relevant at a 0.10 action level, which locate where further records could change a classification and set no action; median revisit interval annual or longer in all 145 cells with one recorded & Ordering may sequence confirmatory sampling for conductivity and TDS (CGWB portal) and fluoride and nitrate (CGWB Year Book); pH not supported; arsenic insufficient support; no pair licenses a probability & Rank prioritisation plus direct confirmation (4); Monitoring baseline required (2); Direct measurement required (1); Targeted direct reference measurement (1) & Local direct measurements permitting recalibration; persistent well, date and sample identity; arsenic method and detection-limit metadata; accepted turbidity and \textit{E.~coli} populations \\ \hline
Kenya, Kwale County (groundwater and springs) & Records above threshold: pH 41 of 140, conductivity 19 of 140, nitrate 5 of 139, fluoride 3 of 139, arsenic 3 of 140; 67 sites sampled in both 2016 campaigns & None; no model was fitted or scored & Local signal plus direct confirmation (5); Monitoring baseline required (3) & Repeat same-site measurement across further seasons of the year; 9 of 24 sites outside the pH range in March were inside it in June \\ \hline
Malawi, laboratory snapshot & Records above threshold: pH 7 of 31, nitrate 0 of 31, turbidity 2 of 22; faecal coliform detectable in 4 of 19; conductivity, fluoride and arsenic not readable as recorded, and TDS on the same mixed scale as conductivity in 10 of 31 records, so no TDS count is given & None; no model was fitted or scored & Local signal plus direct confirmation (4); Data repair required (3); Organism-specific baseline required (2) & Recovered conductivity and TDS scale, fluoride codebook and arsenic units; organism-specific \textit{E.~coli} measurement; wider repeat-site coverage \\ \hline
Malawi, conductivity logger series & 19 sites, 714 site-days over 39 days; 76 site-days above 1,500 \textmu{}S/cm; 2 sites above on every recorded day; 0 changes of state in 695 transitions & None; temporal description only & High-frequency sentinel plus chemical confirmation (1) & Direct chemistry with calibration, fouling, downtime and energy records at the persistent and comparison sites \\ \hline
Bangladesh (published evidence only) & One kit classified 110 of 123 wells correctly against 10 \textmu{}g/L and 113 against the Bangladesh standard of 50 \textmu{}g/L; all 13 errors at 10 \textmu{}g/L were underestimates, and at 50 \textmu{}g/L errors ran in both directions; in tests controlling colour-matching error, two of eight kits were accurate and precise & None; no model transferred & Screening under a stated rule, with reference confirmation of positives and of a sample of negatives; the threshold and the fraction confirmed are the programme's to set & A false-negative rate near the threshold, with its interval, under the programme's rule; operator and invalid-test rates \\ \hline
Ghana (published evidence only) & Modelled cost per test favoured centralised testing in two of three study areas in Ghana and Uganda; in the Ghanaian area it was cheapest in 96\% of simulations; no reviewed field test met the UNICEF target product profile & None; no model transferred & Choose a service architecture on cost per valid result, with the instrument as one component & Network-specific cost per valid result, including invalid and repeat tests, and the sampling frequency at which onsite testing becomes cheaper \\ \hline
\end{longtable}
\begin{flushleft}
1. Counts of records above a threshold are record fractions within each archive, not prevalence or exposure, and a state or Union Territory cell is not a well or a person. The thresholds are those of Materials and methods, applied to every archive; arsenic 0.010 mg/L is 10 \textmu{}g/L.

2. Kwale County and the two Malawi archives are bounded local applications and not national samples. The Kwale counts exclude surface-water samples, and two campaigns three months apart are not a seasonal cycle. Two of the three Kwale fluoride exceedances and two of the 19 conductivity exceedances come from one thermal spring, and the third fluoride exceedance is an unverified value.

3. The Malawi TDS context keeps its recorded action; applying the conductivity unit rule to TDS would return data repair instead.

4. Bangladesh and Ghana had no archive analysed and no model transferred, and they are outside the 26 monitoring contexts. The kit counts are for one kit at the two thresholds stated, the two-of-eight result holds in tests that controlled for colour-matching error, and the costs are modelled.

5. 1,500 \textmu{}S/cm is an operational project convention and not a health guideline. \textit{Ordering} is a model's ranking of sources; ordering that transfers licenses the sequencing of confirmatory sampling, never a probability.

6. No product is named in any setting: the product-level determination is one determination over the 23-product catalogue, applied to every context.

7. \nameref{S12_Table} adds, for each setting, the monitoring decision it poses and where the procedure declines.
\end{flushleft}
}
\linenumbers

\subsection*{India: where further records could change a classification}

The Indian archives resolve to 186 state or Union Territory \texttimes{} parameter cells in 31 states and Union Territories, of which 160 met the support requirement for description. Accepted record counts range from 173,029 for conductivity to 9,927 for arsenic (Table~\ref{table2}).

At an assumed action level of 0.10, 17 of the 160 eligible cells are decision-relevant, in 15 states and Union Territories (Fig~\ref{fig4}A and 4B): 6 fluoride cells, 4 pH, 4 TDS, 2 nitrate and 1 conductivity, and none of the 19 eligible arsenic cells. The count depends on the level: at 0.05 it is 29 cells in 16 states and Union Territories, and at 0.25 it is 4 cells in 4. The intervals treat records as independent although sites are revisited, so the counts are lower bounds. The 17 cells mark where further records could change a state-level classification, not the highest exceedance rates or a ranking of places by hazard; no action in Table~\ref{table1} depends on them, and a cell resolved below the action level is not one where monitoring may stop.

\begin{figure}[!h]
\caption{\textbf{In the Indian archives, 17 of 160 eligible cells are decision-relevant at a 0.10 action level, and fewer cells have the record support for a leave-one-state-out evaluation than for description.} (A) The 17 decision-relevant state or Union Territory \texttimes{} parameter cells, numbered from the widest Wilson 95\% interval to the narrowest, an order of uncertainty and not of severity. Markers are administrative anchors inside each state or Union Territory, not monitoring sites. (B) Exceedance proportion with its Wilson 95\% interval for the same cells; each interval straddles the 0.10 action level, and n is the number of accepted records. The intervals treat records as independent although sites are revisited, so 17 is a lower bound. Decision-relevant is not a severity ranking, a prevalence estimate or a hotspot map. Nitrate uses 50 mg/L as NO\textsubscript{3}; conductivity uses 1,500 \textmu{}S/cm, an operational convention and not a health guideline. (C) Cells per parameter supporting description (grey bar), a leave-one-state-out evaluation (hatched bar) and decision-relevant (diamond): 160, 89 and 17 of 186. The numbers beside each bar give the same three counts for that parameter, in that order. (D) Median revisit interval in the 145 cells with one recorded: none is below 365 days and 48 are exactly 365; the 46 Year Book cells cannot fall below the spacing between editions. TDS, total dissolved solids. Full table: \nameref{S3_Table}. Boundaries: Natural Earth (public domain). Plotted values: \nameref{S9_Data}.}
\label{fig4}
\end{figure}

The other 143 eligible cells are resolved, with intervals wholly above or wholly below the action level. Rajasthan shows the upper side. Of 12,349 conductivity records at 1,641 sites between 2000 and 2024, 6,676 exceeded the 1,500 \textmu{}S/cm project convention (54.1\%, Wilson 95\% interval 53.2\% to 54.9\%), and of 3,131 TDS records at 1,085 sites, 1,627 exceeded 1,000 mg/L (52.0\%, 50.2\% to 53.7\%). No Rajasthan conductivity or TDS record states its analytical method.

In all 145 cells with a recorded median revisit interval, the median is annual or longer; 48 are exactly 365 days and the longest is 1,534 days (Fig~\ref{fig4}D). Of the 145, 46 come from Year Book tables, whose intervals cannot fall below the spacing between editions, and some medians rest on a single interval, so the result describes the archive as held; no action in Table~\ref{table1} prescribes a sampling frequency. Record support would permit a leave-one-state-out evaluation in 89 of the 186 cells, including 9 of the 17 decision-relevant cells, and in none of the 31 arsenic cells. The arsenic record fraction of 0.0031 is not interpretable as a monitoring signal until units, method and detection limits are recovered.

\subsection*{What the model evidence permits}

How far the model generalises depends on the validation design (Fig~\ref{fig3}, Table~\ref{table3}). Across seven tasks from the multinational fitting archive, mean ROC-AUC on the test partition was 0.9396 under random-row splitting, 0.8688 under station-grouped splitting and 0.7696 under 3\textdegree{} block splitting, an optimism of 0.1701 computed from unrounded means. The loss from random-row to block splitting ranges from 0.0768 for fluoride to 0.2927 for nitrate. In forward-chained yearly folds, Model B exceeded Model A on each of the six parameters with a paired comparison, by 0.0215 to 0.0643 in ROC-AUC, an advantage that needs earlier measurements at the same site.

Of the 13 archive \texttimes{} parameter pairs evaluated for transfer (Fig~\ref{fig5}), ordering transferred in the Indian archives for conductivity and TDS in the CGWB portal and for fluoride and nitrate in the CGWB Year Book. It was not supported for pH in either archive in which pH was tested (threshold ROC-AUC 0.4870 in the CGWB portal and 0.4886 in WQP). Arsenic in the CGWB portal had insufficient support, with one evaluable test on 119 scored observations at 116 sites, and its ordering was inverted: ROC-AUC 0.3170 (95\% interval 0.2043 to 0.4450). TDS was the only parameter supported in both archives in which it was tested (threshold ROC-AUC 0.7081 in the CGWB portal and 0.7406 in WQP). The nitrate ordering was verified at 10 mg/L as N, and its use against 50 mg/L as NO\textsubscript{3} assumes that the ordering does not depend on that difference. Scored observations are fewer than accepted records for every Indian pair; the site exclusion accounts for part of the difference, and the rest is not recorded.

\begin{figure}[!h]
\caption{\textbf{Transfer to external archives: rank ordering was supported for 5 of 13 archive \texttimes{} parameter pairs, with calibration slopes of 0.2383 to 0.4017 on those five.} (A) The verdict for each archive \texttimes{} parameter pair, with threshold ROC-AUC and the number of external observations; 13 of 24 pairs were evaluated, and the 11 untested pairs are unknown evidence. The verdict comes from three interval tests, on site-level rank correlation, on rank correlation within 1\textdegree{} spatial blocks and on lift in the top predicted decile, not from ROC-AUC; insufficient support is never merged with not supported. In B--E the marker shape repeats the verdict: circle, rank transfer supported; square, partial or scale-dependent; triangle, insufficient support; cross, not supported. (B) Threshold ROC-AUC with its 95\% site-clustered bootstrap interval (2,000 replicates; 1,998 for Water Quality Portal (WQP) nitrate); the CGWB portal pH interval, 0.4784 to 0.4964, lies wholly below 0.5. (C) The Spearman rank correlation between predicted risk and observed severity among detected records, with its site-clustered 95\% interval; it is not one of the verdict tests. Severity is log\textsubscript{10} concentration, or for pH the distance from the centre of the 6.5--8.5 band. (D) Calibration slope, below 1 on all 13 pairs and at or below 0 on 4. (E) Brier skill against a reference that knows the external base rate, negative on 11 of 13 pairs and on all six Indian pairs. Calibration-in-the-large was not assessed and no local recalibration was run, so no pair licenses a model probability. The fitting archive is predominantly surface water and both external archives are groundwater; nitrate is scored at 10 mg/L as N. CGWB, Central Ground Water Board; TDS, total dissolved solids. Plotted values: \nameref{S9_Data}.}
\label{fig5}
\end{figure}

Probability did not transfer with ordering. The calibration slope was below 1 on all 13 evaluated pairs, between 0.2383 and 0.4017 on the five pairs whose ordering transferred, and at or below 0 on 4. Within the fitting archive the block-design slopes were already below 1 for all seven tasks, between 0.1925 and 0.8922. Brier skill was negative on all six Indian pairs and on 11 of the 13 overall; the two non-negative values, 0.0470 for \textit{E.~coli} and 0.0169 for arsenic in WQP, belong to pairs without supported rank transfer (Fig~\ref{fig5}D and 5E). The conductivity base rate was 0.0638 in the transfer training data and 0.2364 among the scored CGWB portal observations, against 0.2261 among all accepted conductivity records. Calibration-in-the-large was not assessed and local recalibration was not run, so no transferred score can be read as a local exceedance probability.

At a fixed 0.5 cut-point most exceeding observations were not flagged; for conductivity in the CGWB portal, sensitivity was 0.0787 (\nameref{S1_Text} gives the counts).

For India the rule returned four actions (Table~\ref{table1}). Conductivity, TDS, nitrate and fluoride returned \textit{rank prioritisation plus direct confirmation}: ordering may set the order in which sources are confirmed, but every flagged source still needs a direct measurement, low-ranked sources stay under routine measurement, and no ranking is released with this article. pH returned \textit{direct measurement required}, and arsenic \textit{targeted direct reference measurement}, pending method and detection-limit metadata, with no basis set for choosing targets and laboratory access and cost not assessed. Turbidity and \textit{E.~coli}, with no accepted Indian measurement population in the archives analysed, returned \textit{monitoring baseline required}; CGWB lists microbial contaminants among additional parameters to be monitored~\cite{CGWB-SOP-GWQDataAnalysis}, so the evidence may lie with drinking-water agencies rather than in a new survey. Tables~\ref{table1} and~\ref{table4} give the evidence that would change each action.

\subsection*{Kenya: repeat measurement in Kwale County}

The Kwale archive describes groundwater and springs sampled in Kwale County and supports no statement about Kenya as a whole. It records 81 sampling locations, counted as distinct coordinates across 71 named localities, 73 of which remain once surface-water samples are excluded, visited in March and June 2016, with 139 or 140 records per parameter. Measured pH lay outside 6.5--8.5 in 41 of 140 records (29.3\%, Wilson 95\% interval 22.4\% to 37.3\%), all of them below 6.5, and conductivity exceeded the 1,500 \textmu{}S/cm convention in 19 of 140 (13.6\%, 8.9\% to 20.2\%). Nitrate exceeded 50 mg/L as NO\textsubscript{3} in 5 of 139 (3.6\%, 1.5\% to 8.1\%), fluoride exceeded 1.5 mg/L in 3 of 139 (2.2\%, 0.7\% to 6.2\%) and arsenic exceeded 10 \textmu{}g/L in 3 of 140 (2.1\%, 0.7\% to 6.1\%). Most locations were sampled twice, so these intervals understate the uncertainty.

Two of the three fluoride exceedances, and two of the 19 conductivity exceedances, come from one thermal spring with a recorded temperature of about 59 \textdegree{}C, and the third fluoride exceedance and one June conductivity crossing are unverified values (\nameref{S1_Text}). The counts include these records as the archive holds them.

Of the 67 sites sampled in both campaigns (66 with nitrate and fluoride results in both), 15 of the 24 with pH outside the range in March were still outside it in June and 2 more crossed; 8 of the 10 above the conductivity convention in March remained above and 1 more crossed; and for nitrate, fluoride and arsenic one site was above the threshold in both campaigns, with 2, 1 and 1 crossing only in June (\nameref{S11_Table}); the fifth nitrate exceedance is at a location sampled once. The source labels the campaigns dry-season and wet-season, but one campaign in each season of a single year does not describe a seasonal cycle.

No model was used for Kwale. For pH, conductivity, nitrate, fluoride and arsenic, the rule returned \textit{local signal plus direct confirmation} from the data status alone: repeat measurement at the same sites across further seasons, with threshold classifications confirmed by a reference method the rule does not name (Table~\ref{table1}). Contexts that share the archive and the action could be met by one sampling round. The action sets what to measure next; it does not postpone confirmation at, or a programme's response to, a source recorded in 2016 above a WHO health-based guideline value. TDS, turbidity and \textit{E.~coli}, for which the archive holds no records, returned \textit{monitoring baseline required}.

\subsection*{Malawi: a laboratory snapshot and a logger series}

The Malawi laboratory snapshot holds one record per water point for 32 water points sampled between December 2017 and July 2019, and it is not a national sample. Measured pH lay outside 6.5--8.5 in 7 of 31 records (22.6\%, Wilson 95\% interval 11.4\% to 39.8\%), all between 6.24 and 6.42, and the archive flags each as within the Malawian standard it applies. No nitrate record of 31 exceeded 50 mg/L as NO\textsubscript{3}, which leaves an upper Wilson bound of 11.0\%. Turbidity exceeded the 1 NTU rule in 2 of 22 records (9.1\%, 2.5\% to 27.8\%), both flagged as within the Malawian standard. Faecal coliform was detectable in 4 of 19 samples, all collected in February 2019 and none with a recorded analysis date. Faecal coliform is not \textit{E.~coli}, but WHO treats thermotolerant coliforms as an acceptable alternative indicator of faecal pollution~\cite{WHO2022GDWQ}, so detection at those four water points would ordinarily warrant inspection and resampling whatever the organism-specific evidence gap.

Four parameters could not be read as recorded. Conductivity failed a unit-consistency check, TDS is recorded on the same mixed scales in 10 of 31 records, so no TDS count is reported, fluoride held only a placeholder code, and arsenic had no documented unit. Conductivity, fluoride and arsenic returned \textit{data repair required}, to be made from recovered records and not by assumption; only the originating laboratory or the data publisher can make it, and the rule has no fallback if the records cannot be recovered. The laboratory pH, TDS, nitrate and turbidity contexts returned \textit{local signal plus direct confirmation}, although under the conductivity rule the TDS context would return data repair, and for nitrate the action amounts to wider coverage. \textit{E.~coli}, with no organism-specific measurement, and the faecal coliform context both returned \textit{organism-specific baseline required} (Table~\ref{table1}).

The logger extract holds conductivity from 19 sites in 9 districts, reduced to 714 daily site medians between 31 December 2024 and 7 February 2025, a window of 39 days within Malawi's wet season from November to April~\cite{BGS-Earthwise-Malawi}. On 76 of the 714 site-days, conductivity exceeded the 1,500 \textmu{}S/cm convention: two sites, both in Chikwawa District, exceeded it on every recorded day and no other site on any day, so across 695 successive site-day transitions no site changed state. The two sites have median conductivities of about 5,690 and 2,160 \textmu{}S/cm, and the next highest site median is about 1,360 \textmu{}S/cm. The logger does not identify the cause, since conductivity is not specific to any contaminant, does not record whether values are temperature-compensated, and says nothing about the dry season. The rule returned \textit{high-frequency sentinel plus chemical confirmation}, pairing the two persistent sites, and comparison sites, with direct chemistry and with instrument calibration, fouling, downtime and energy records; the analytes, comparison sites, chemistry frequency and logging period are not set by this study.

\subsection*{Bangladesh and Ghana: decisions from published evidence}

In Bangladesh and Ghana no archive was analysed and no model was transferred, and neither setting is among the 26 contexts. The Bangladesh decision is whether a field kit can classify a well near the arsenic threshold accurately enough to leave it in service. In a blind evaluation of 123 wells against laboratory measurement, with 65 tested twice, one kit classified 110 wells correctly against the WHO guideline of 10 \textmu{}g/L and 113 against the Bangladesh standard of 50 \textmu{}g/L~\cite{George2012}. All 13 of its misclassifications against 10 \textmu{}g/L were underestimates, in wells between 10 and 27 \textmu{}g/L~\cite{George2012}. Against 50 \textmu{}g/L its over- and underestimates were evenly distributed, and the evaluation omitted the kit's hydrogen-sulphide oxidation step~\cite{George2012}. A comparison of eight commercial kits with hydride-generation atomic absorption spectrometry found two accurate and precise in tests that controlled for colour-matching error, four accurate or precise but not both, and two neither~\cite{Reddy2020}.

Applied to that evidence, the procedure permits screening under a stated decision rule with confirmation of screen positives and of a sample of negatives by a reference method. The threshold and the fraction of negatives confirmed are not set by this study; at 10 \textmu{}g/L the recorded errors all missed contaminated wells, which is the case for confirming negatives, whereas at 50 \textmu{}g/L they ran in both directions. Kit and laboratory costs in Bangladesh were not assessed. The evidence that would change the action is a false-negative rate near the threshold, with its interval, under the rule a programme uses, together with operator and invalid-test rates.

The Ghana decision is how to organise \textit{E.~coli} testing so that frequency, transport, incubation and staffing remain affordable. Modelled costs per membrane-filtration test for four service approaches, in three study areas in Ghana and Uganda, favoured centralised testing in two of the three~\cite{Trimmer2024}. In the Ghanaian study area, centralised testing was cheapest in 96\% of simulations, at a median of USD 10 per test~\cite{Trimmer2024}. The same study reports the WHO recommendation that samples not be stored for more than 6 h before processing, which links distance to the choice of service, and publishes its cost model for others to run~\cite{Trimmer2024}. A review of commercial field tests found none that met the UNICEF target product profile, with the need for a power source the largest limitation and 18 to 24 h of incubation still required~\cite{Pichel2023}.

The procedure therefore treats the instrument as one component of a service architecture chosen on cost per valid result, which counts invalid and repeat tests; a valid microbial result is one processed within the holding time, incubated as specified and accompanied by passing controls. Until that is available, modelled cost per test by service approach, bounded to its study areas and assumptions, is the basis a programme has. Neither costing study states a price year, although one converted costs collected in December 2014 and January 2015 at 1 January 2015 exchange rates~\cite{Delaire2017}, and none of the 23 products in the documentary catalogue documents a per-test cost with its basis, currency and date. Neither setting's evidence supports naming a kit or a product.

\subsection*{Product documentation across the settings}

Across the settings, the technology question meets the same documentary record. The exact-product grid holds 23 commercial products against 54 criteria, 1,242 cells (Fig~\ref{fig6}A). Of these, 891 (71.7\%) are not reported in any located document, 195 (15.7\%) are document-sourced, 9 of them from regulatory certification, and 106 were not researched, all for two products (Fig~\ref{fig6}A gives the remaining classes). Not reported means that the located documentation is silent, not that a product was measured and found wanting. The 9 regulatory cells all belong to one product; they are independent of the manufacturer but are documentation, not device-versus-reference testing.

\begin{figure}[!h]
\caption{\textbf{The documentary evidence for 23 named products, and the documentation threshold for the 11 active criteria.} (A) All 1,242 cells of 23 products \texttimes{} 54 criteria, in five families, by evidence class: 891 not reported in any located document, 106 not researched (all for two products), 40 conflicting, adjacent-product or range-only, 10 not applicable, and 195 document-sourced, 9 of them from regulatory certification for one product, which is documentation and not device-versus-reference testing. Products are shown by identifier; identities and sources are in \nameref{S4_Table}. MTBF, mean time between failures; QC, quality control; CV, coefficient of variation; IP, ingress protection. A not-reported cell is unknown evidence about a product, not evidence about its performance. (B) The 11 active criteria against the documentation threshold of 12 of 23 products: 1 clears it on strict primary fields and 4 on any admissible field. No product documents all 11 and the most any documents is 9, so the procedure returns a product-level abstention, a statement about the evidence base and not a ranking. Plotted values: \nameref{S9_Data}.}
\label{fig6}
\end{figure}

Under the two counting rules set out in Materials and methods, 1 of the 11 active criteria meets the documentation threshold on the strict field and 4 on any admissible field (Fig~\ref{fig6}B). No product documents all 11 active criteria, including the 21 that were researched. On any admissible field the most any product documents is 9, reached by 1 product, with 4 products documenting 8 and 4 documenting 7; on the strict field the most is 7. The outcome therefore rests on requiring all 11: a rule requiring 9 would admit one product. For 19 of the 54 criteria, no located document covers any product.

Every one of the 26 contexts, each a country \texttimes{} source archive \texttimes{} parameter combination, returned an action, and the 26 took eight distinct actions (Fig~\ref{fig1}B, Table~\ref{table1}). The most frequent were \textit{local signal plus direct confirmation}, in 9 contexts, and \textit{monitoring baseline required}, in 5. Six of the 26 contexts hold no records, and their actions follow from that absence. In all 26, the procedure named no product: one determination over the 23-product catalogue applied to every context, not 26 separate determinations and not evidence about the contexts.

Among technology archetypes, a population separate from the named products, a single archetype led under every cost mapping in 5 of 7 weighting scenarios, and the leader changed with the cost mapping in the other two (\nameref{S7_Table}); because the comparison was not screened for analytes, a leading archetype is preference context under stated weights, not a product or the class of instrument to use.

\section*{Discussion}

\subsection*{What the settings show}

In India the evidence supports a sequence rather than a map: ordering could decide which conductivity, TDS, nitrate or fluoride sources to confirm first, not which may be left untested, and resolved cells such as Rajasthan call for contaminant-specific confirmation rather than more records. Because calibration was poor on every transfer pair, no Indian programme could read a transferred score as the probability that a source exceeds its guideline, and because the transfer also crossed water-body type, its supported pairs cannot separate a change of geography from a change of water type. We offer this as a reading of the results, not as a tested operational gain.

In Kwale County and Malawi no model output entered any decision, and the constraint was the record itself. In Kwale the binding step is returning to the same sources, since pH classifications changed between campaigns at 11 of the 67 sites sampled twice, and the fluoride signal rests on a thermal spring and an unverified value that a practitioner would check before acting. In the Malawi laboratory snapshot, three parameters cannot be read until units, a missing-value codebook and, for arsenic, detection limits are recovered; these gaps describe the records and say nothing about the water. The logger series shows persistence at two sites without identifying its cause. Both sites are in Chikwawa District, and groundwater salinity too high for potable use has been documented in Malawi's lower Shire Valley~\cite{BGS-Earthwise-Malawi}; whether that explains these sites was not assessed. The snapshot's gaps are in line with reports that groundwater monitoring data in Malawi are persistently scarce, a scarcity made worse by existing monitoring networks that are not operating~\cite{Robertson2026-PLOSW0000593}.

In Bangladesh and Ghana the constraint lay in the decision rule and the service design rather than in any model. In hypothetical well-switching scenarios for Araihazar, Bangladesh, relying on kit rather than laboratory data lost only a modest part of the reduction in average arsenic exposure, set against the logistical and financial challenge of supplying laboratory data alone~\cite{Jameel2021}. The same analysis found the 50 \textmu{}g/L threshold used in Bangladesh, rather than the WHO guideline of 10 \textmu{}g/L, close to optimal for average exposure reduction~\cite{Jameel2021}. A programme could accept kit screening on that basis, but the residual error it will tolerate is the programme's choice and not a result of this study. For Ghana, both costing studies cited here, one across sub-Saharan Africa and one covering study areas in Ghana and Uganda, express cost per test with quality-control samples added as a markup~\cite{Delaire2017,Trimmer2024}, whereas a service design needs cost per valid result, counting invalid and repeat tests.

\subsection*{Relation to previous work}

Mean discrimination fell from 0.9396 to 0.7696 between two validation designs applied to the same models and data, whereas a review of machine-learning models for United States drinking water found tenfold cross-validation or a randomly chosen test set the most common basis for reported performance~\cite{Hu2023}. In a forest biomass mapping study, non-spatial validation suggested that a model predicted more than half of the variation, whereas spatial validation found almost no predictive power~\cite{Ploton2020}. A methodological source on spatial dependence records the counter-argument that spatial cross-validation is not the right way to evaluate map accuracy~\cite{MeyerPebesma2022}, so the single spatial-block split is used here to address a generalisation claim, not to certify accuracy away from sampled locations. Whether the national hazard models for India cited in the Introduction carry such optimism was not tested here.

The calibration result also has precedents. Estimated risks can be unreliable even when discrimination is good, for example overestimating risk when a model is used in a setting where incidence is lower~\cite{VanCalster2019}. Calibration-in-the-large, the calibration slope and discrimination have been proposed as distinct measures of model performance, with problems of calibration-in-the-large often found at external validation~\cite{SteyerbergVergouwe2014}, and moderate calibration has been proved to guarantee decision-making that does no harm~\cite{VanCalster2016}. No transferred score in this study was shown to be moderately calibrated, so none carries that guarantee, and because calibration-in-the-large was not assessed, the study shows that transferred probabilities were miscalibrated without describing every form of that miscalibration.

Evidence from outside the product catalogue points the same way as the documentary shortfall. Besides the operator-dependent kit verification cited in the Introduction, a laboratory evaluation of a marketed low-cost sensor found no correlation between its score and the presence of \textit{E.~coli}~\cite{Purvis2026-PLOSW0000442}. Both evaluated one product at a stated date, and no conclusion about any product is drawn from them. Most of 72 institutions assessed across ten sub-Saharan African countries did not achieve the microbial testing levels set by applicable standards or WHO guidelines~\cite{Peletz2016}, and drinking-water surveillance often falls short of its required frequency~\cite{GLAAS2025}. Threshold-near performance, operator dependence and achievable testing frequency are what product documentation would need to report for a programme to rely on a named instrument.

Abstention here differs from its usual machine-learning form, in which a model abstains where it is likely to make a mistake~\cite{Hendrickx2024}. The procedure abstains where the evidence needed for a decision is absent, so its abstention describes the evidence base rather than predictive error. Stochastic multicriteria acceptability analysis aids decisions in problems with inaccurate, uncertain or missing information~\cite{TervonenLahdelma2007}, and the archetype comparison draws on it, but no sampling of weights can supply a criterion value that no located document holds. We therefore read the documentary shortfall as the result and the product-level abstention as the output of the sufficiency rule applied to it.

\subsection*{Limitations}

The model evidence has five limits. The external archives are not demonstrably independent of the fitting archive by record, and the transfer crosses water-body type as well as geography. The optimism of 0.1701 is specific to two designs on seven tasks, and the transferred models were separate fits never evaluated under those designs, so the optimism cannot be carried to the transferred scores. The transfer verdicts rest on tests whose statistics were not retained, so they carry no stated sampling uncertainty. Calibration was assessed without calibration-in-the-large, local recalibration was not run because no prediction-level file was retained, and no fitted object is recorded for the validation or transfer models. The runs behind Table~\ref{table3}, Fig~\ref{fig5} and \nameref{S5_Text} were not re-run for this article, so their values are verified against retained result exports, and a result-table verification is not a full training reproduction.

The archive evidence has its own limits. Record fractions describe what was sampled, not the population served. Wilson intervals treat records as independent although sites are revisited, so the 17 decision-relevant cells are a lower bound, and the intervals do not remove sampling-design bias. CGWB's current design samples every quality station once every five years before the monsoon and at least a quarter of the stations before and after the monsoon every year~\cite{CGWB-SOP-GWQDataAnalysis}, so the revisit intervals reported here describe the archive as held rather than the programme as now designed. Acceptance of the local archives applied no charge-balance or conductivity--TDS consistency check, and no records from other Indian or Kenyan custodians, microbial records included, were used. The Kwale and Malawi archives are small and local, the Kwale location count depends on how coordinates were recorded, and the logger extract covers 39 days. The Bangladesh and Ghana studies were not found by a systematic search, so other studies could change the evidence those applications rest on.

The decision evidence is limited in five further ways. The evidence states behind the six Indian contexts with accepted records were assigned by the authors from the transfer verdicts and record support, and no dated protocol records when the decision rules were fixed. No record states how the 23 products were chosen or which sources were searched and when, so the product counts describe the documentary record as located; a criterion no located document reports may still have been measured, and two products were not researched at all. The abstention depends on requiring all 11 active criteria, and a rule requiring 9 would admit one product. Among the records assembled for the 23 products, paired device-versus-reference field evidence was located only for arsenic, and only one of the three recorded comparisons has recoverable publication provenance. No practitioner weights were elicited, no monitoring programme reviewed the returned actions for feasibility or cost, the procedure was not compared with an alternative procedure, and no one outside this project has applied its rules, so whether other analysts would return the same actions is untested.

\subsection*{What would test or change these results}

Two observations would test the procedure. It would gain support if the same rules, applied to a catalogue with fuller documentation, admitted a named product, since that would show that the abstention reported here reflects this evidence base rather than rules no product could meet. It would be shown to be underspecified if two analysts applying the stated rules to the same archive returned different actions. The evidence that would change the actions reported here is specific: local measurements that allow model recalibration (\nameref{S8_Table}, priority D4), persistent site and well identity (D2), paired chemistry at sentinel sites (D3), threshold-near kit performance under a stated rule (A1), and invalid and repeat tests recorded in a deployed network, from which cost per valid result follows (C1). The measurements a recalibration needs are themselves direct measurements, so its value to a programme depends on the model then reducing later measurement, which was not assessed. \nameref{S8_Table} sets out 16 research priorities in four families; it does not cover a comparison of service architectures on cost per valid result, the local conductivity--TDS relation, or the two-analyst test proposed above.

\section*{Conclusion}

Applied to India, Kwale County in Kenya, Malawi, Bangladesh and Ghana, the procedure returned what the evidence could support in each setting. In India, ordering transferred for four archive \texttimes{} parameter pairs (conductivity and TDS in the CGWB portal, fluoride and nitrate in the CGWB Year Book) and could sequence confirmatory sampling; at an assumed 0.10 action level, 17 of 160 eligible cells were decision-relevant, a lower bound; and no transferred score could be read as a probability. In the bounded local archives of Kwale and Malawi, where no model was used, the actions were repeat measurement, baselines, data repair and paired chemistry. In Bangladesh and Ghana, published evidence supported screening with confirmation and a service design chosen on cost per valid result, which neither costing study reports. In all 26 monitoring contexts, the procedure named no product.

That abstention is an output, and the result behind it is the documentary record: most product documentation cells are not reported in any located document, and no product documents all 11 active criteria. The evidence that would change these actions is specific: measurements allowing model recalibration, persistent site identity, paired chemistry, threshold-near kit performance and cost per valid result. Until it exists, a programme could use ordering to sequence, not to reduce, confirmatory sampling where it transferred, and direct measurement to make each decision; this is a reading of the results, not a tested operational gain.

\section*{Supporting information}

\paragraph*{S1 Text.}
\label{S1_Text}
\textbf{Harmonisation, model configuration and decision rules.} Harmonisation, unit and chemical-basis rules, with the Indian standard's limits and the arsenic unit remedies; the accepted Indian populations and their record counts; the data states recorded in place of removed records; how the 26 monitoring contexts were generated; the versions of the WQP extract and the Malawi snapshot; the validation-design and transfer-fit configurations, interval replicates and confusion counts at the 0.5 cut-point; the decision-rule branches, the 11 active criteria and the dated records of when rules were set; the Kwale values that are not verified; the training-set sizes of the transfer tasks; the Year Book source documents, with their URLs, Internet Archive snapshot identifiers and SHA-256 digests; and the 62 National Water Data Portal resource files, with their byte counts.

\paragraph*{S2 Table.}
\label{S2_Table}
\textbf{The 26 monitoring contexts in full.} Every field of the context matrix, of which Table 1 carries the decision columns, and the definitions of the eight actions: one row per country \texttimes{} source archive \texttimes{} parameter combination. The product-level decision is one determination applied to every row.

\paragraph*{S3 Table.}
\label{S3_Table}
\textbf{Support, cadence and exceedance for 186 Indian state or Union Territory \texttimes{} parameter \texttimes{}
source cells.} Table 2 and Fig 4 summarise it. Exceedance proportions are of accepted records, with Wilson 95\% intervals; a proportion of records is not a proportion of people.

\paragraph*{S4 Table.}
\label{S4_Table}
\textbf{The 23 \texttimes{} 54 product documentation grid.} Product identities (P01--P23 in Fig 6) and, for each cell, the evidence class, value, unit, source document and locus. A not-reported cell is unknown evidence about a product, not evidence that the product performs poorly.

\paragraph*{S5 Text.}
\label{S5_Text}
\textbf{Mechanistic ablation.} Rung definitions, configuration, layer sources, and the discrimination of each rung with its block-bootstrap interval.

\paragraph*{S6 Fig.}
\label{S6_Fig}
\textbf{An illustrative conductivity prediction surface.} A 60 km window centred on 24.0\textdegree{}N, 87.5\textdegree{}E, chosen by the number of monitoring sites with complete coverage of the 16 mechanistic features, without reference to any conductivity value or model output. (A) The output of the coordinate-free full-stack rung (A5) refitted under the frozen configuration and without tuning, with every three-degree block intersecting the window excluded from training --- two blocks, because the window straddles 24\textdegree{}N --- so no training observation lies inside the area shown: 148,029 training rows, 7,421 test observations at 1,613 test sites, the 1,500 \textmu{}S/cm conductivity convention (an operational screening value, not a health guideline) and a training prevalence of 0.2292. The display refit carries no calibration result, so its scale is read as an ordering, not a probability. (B) Soil pH (water), 0--5 cm, from SoilGrids 2.0. (C) Population density in 2020, people per km\textsuperscript{2} on a logarithmic scale, from WorldPop. Grey cells hold no value. The land-cover, aquifer, clay and organic-carbon layers of the stack are not drawn, and no monitoring site is shown. The surface is a spatial illustration, and a modelled cell is not a well, tap, household or person.

\paragraph*{S7 Table.}
\label{S7_Table}
\textbf{Rank acceptability of technology archetypes in seven weighting scenarios.} For each scenario, the archetype that leads among the ten compared, its mean probability of ranking first and the range of that probability across 13 cost mappings, its expected rank, and whether the same archetype leads under every mapping. The archetypes are a separate population from the named products; no row is a product recommendation, and the product-level outcome is abstention in every context.

\paragraph*{S8 Table.}
\label{S8_Table}
\textbf{Sixteen research priorities in four families.} Each priority with its family, priority level, proposed study design, measurement endpoint and the decision it would improve.

\paragraph*{S9 Data.}
\label{S9_Data}
\textbf{The values plotted in Figs 1--6 and S6 Fig.} One CSV file per figure, written by the figure builders from the numbers they drew.

\paragraph*{S10 Text.}
\label{S10_Text}
\textbf{Record of artificial-intelligence assistance.} Each activity, its inputs and outputs, and the verification applied to it.

\paragraph*{S11 Table.}
\label{S11_Table}
\textbf{Record fractions for the Kenyan and Malawian archives, Kwale paired campaigns and the Malawi
logger summary.} Eleven record fractions with Wilson 95\% intervals, which treat records as independent and so understate uncertainty; threshold states at the Kwale sites sampled in both 2016 campaigns, with surface-water samples excluded; and the summary of the 39-day logger extract, all recomputed from the raw openwashdata archives. Record fractions are not prevalence, and neither country's archives are national samples. The Malawi TDS fraction is not interpretable, because TDS is recorded on mixed scales in 10 of 31 records, and the article reports no TDS count.

\paragraph*{S12 Table.}
\label{S12_Table}
\textbf{The setting synthesis in full.} For each application setting, the monitoring decision it poses, the evidence held, what a model may do, the action returned, the evidence that would change it and where the procedure declines. Table 4 carries five of these seven columns, and its notes apply here.

\section*{Acknowledgments}

We acknowledge the originators of the water-quality records held in GEMStat and the UNEP GEMS/Water Programme, whose database supplied the records used to fit the models. Indian groundwater-quality data were produced by the Central Ground Water Board and obtained through the National Water Data Portal of the National Water Informatics Centre, Ministry of Jal Shakti, and from the Board's Ground Water Quality Data documents. United States records were obtained from the Water Quality Portal of the National Water Quality Monitoring Council, the U.S. Geological Survey and the U.S. Environmental Protection Agency. The Kenyan and Malawian local archives are openwashdata packages. Soil properties are SoilGrids 2.0 data obtained through ISRIC -- World Soil Information; land cover is ESA WorldCover 10 m 2021 v200; and population density is from WorldPop.

\textbf{[AUTHOR DECISION B4: whether to acknowledge supervision or other help, stating the contribution, and only with that person's permission. Anyone who meets the authorship criteria belongs in the author list instead.]}

\nolinenumbers

% References: generated with plos2025.bst from references.bib and embedded here, so this file is self-contained.
\begin{thebibliography}{10}

\bibitem{JMP2025}
{WHO/UNICEF Joint Monitoring Programme for Water Supply, Sanitation and
  Hygiene}.
\newblock {Progress on household drinking water, sanitation and hygiene
  2000--2024: special focus on inequalities}.
\newblock Geneva: World Health Organization and United Nations Children's Fund;
  2025.
\newblock ISBN 978-92-4-011514-9.

\bibitem{GLAAS2025}
{World Health Organization}, {United Nations Children's Fund}.
\newblock {State of systems for drinking-water, sanitation and hygiene: global
  update 2025}.
\newblock Geneva: World Health Organization; 2025.
\newblock ISBN 978-92-4-011898-0.

\bibitem{Bain2014}
{Bain} R, {Cronk} R, {Wright} J, {Yang} H, {Slaymaker} T, {Bartram} J.
\newblock {Fecal contamination of drinking-water in low- and middle-income
  countries: a systematic review and meta-analysis}.
\newblock PLoS Med. 2014;11(5):e1001644.
\newblock \href {http://dx.doi.org/10.1371/journal.pmed.1001644}
  {doi:10.1371/journal.pmed.1001644}.

\bibitem{CGWB-SOP-GWQDataAnalysis}
{Central Ground Water Board}.
\newblock {Standard operating procedure on ground water quality data analysis:
  step by step guidelines for ground water quality data analysis}.
\newblock Faridabad: Central Ground Water Board, Ministry of Jal Shakti; [date
  unknown].
\newblock Available from:
  \url{https://cgwb.gov.in/cgwbpnm/public/uploads/documents/17333891661542642909file.pdf}
  [cited 2026 Sep 15].

\bibitem{Podgorski2020}
{Podgorski} J, {Berg} M.
\newblock {Global threat of arsenic in groundwater}.
\newblock Science. 2020;368(6493):845-50.
\newblock \href {http://dx.doi.org/10.1126/science.aba1510}
  {doi:10.1126/science.aba1510}.

\bibitem{Podgorski2018}
{Podgorski} JE, {Labhasetwar} P, {Saha} D, {Berg} M.
\newblock {Prediction modeling and mapping of groundwater fluoride
  contamination throughout India}.
\newblock Environ Sci Technol. 2018;52(17):9889-98.
\newblock \href {http://dx.doi.org/10.1021/acs.est.8b01679}
  {doi:10.1021/acs.est.8b01679}.

\bibitem{Podgorski2020-IJERPH-India}
{Podgorski} J, {Wu} R, {Chakravorty} B, {Polya} DA.
\newblock {Groundwater arsenic distribution in India by machine learning
  geospatial modeling}.
\newblock Int J Environ Res Public Health. 2020;17(19):7119.
\newblock \href {http://dx.doi.org/10.3390/ijerph17197119}
  {doi:10.3390/ijerph17197119}.

\bibitem{Meyer2019}
{Meyer} H, {Reudenbach} C, {W{\"o}llauer} S, {Nauss} T. {Importance of spatial
  predictor variable selection in machine learning applications -- moving from
  data reproduction to spatial prediction} [Preprint]. arXiv:1908.07805v1;
  2019.
\newblock \href {http://dx.doi.org/10.48550/arXiv.1908.07805}
  {doi:10.48550/arXiv.1908.07805}.

\bibitem{MeyerPebesma2022}
{Meyer} H, {Pebesma} E.
\newblock {Machine learning-based global maps of ecological variables and the
  challenge of assessing them}.
\newblock Nat Commun. 2022;13(1):2208.
\newblock \href {http://dx.doi.org/10.1038/s41467-022-29838-9}
  {doi:10.1038/s41467-022-29838-9}.

\bibitem{NiculescuMizil2005}
{Niculescu-Mizil} A, {Caruana} R.
\newblock {Predicting good probabilities with supervised learning}.
\newblock In: Proceedings of the 22nd International Conference on Machine
  Learning (ICML '05). ACM Press; 2005. p. 625-32.
\newblock \href {http://dx.doi.org/10.1145/1102351.1102430}
  {doi:10.1145/1102351.1102430}.

\bibitem{VanCalster2019}
{Van Calster} B, {McLernon} DJ, {van Smeden} M, {Wynants} L, {Steyerberg} EW;
  {Topic Group `Evaluating diagnostic tests and prediction models' of the
  STRATOS initiative}.
\newblock {Calibration: the Achilles heel of predictive analytics}.
\newblock BMC Med. 2019;17(1):230.
\newblock \href {http://dx.doi.org/10.1186/s12916-019-1466-7}
  {doi:10.1186/s12916-019-1466-7}.

\bibitem{Andres2018}
{Andres} L, {Boateng} K, {Borja-Vega} C, {Thomas} E.
\newblock {A review of in-situ and remote sensing technologies to monitor water
  and sanitation interventions}.
\newblock Water (Basel). 2018;10(6):756.
\newblock \href {http://dx.doi.org/10.3390/w10060756} {doi:10.3390/w10060756}.

\bibitem{EPA-ETV-QuickLowRangeII-2003}
{US Environmental Protection Agency Environmental Technology Verification
  Program; Battelle Advanced Monitoring Systems Center}.
\newblock {ETV joint verification statement: arsenic test kit, Quick Low Range
  II (Industrial Test Systems, Inc.)}.
\newblock Washington (DC): US Environmental Protection Agency; 2003 Sep.
\newblock Available from:
  \url{https://archive.epa.gov/nrmrl/archive-etv/web/pdf/01_vs_lowrangell.pdf}
  [cited 2026 Sep 15].

\bibitem{Delaire2017}
{Delaire} C, {Peletz} R, {Kumpel} E, {Kisiangani} J, {Bain} R, {Khush} R.
\newblock {How much will it cost to monitor microbial drinking water quality in
  sub-Saharan Africa?}
\newblock Environ Sci Technol. 2017;51(11):5869-78.
\newblock \href {http://dx.doi.org/10.1021/acs.est.6b06442}
  {doi:10.1021/acs.est.6b06442}.

\bibitem{Schuwirth2018}
{Schuwirth} N, {Honti} M, {Logar} I, {Stamm} C.
\newblock {Multi-criteria decision analysis for integrated water quality
  assessment and management support}.
\newblock Water Res X. 2018;1:100010.
\newblock \href {http://dx.doi.org/10.1016/j.wroa.2018.100010}
  {doi:10.1016/j.wroa.2018.100010}.

\bibitem{Broekhuizen2015}
{Broekhuizen} H, {Groothuis-Oudshoorn} CGM, {van Til} JA, {Hummel} JM,
  {IJzerman} MJ.
\newblock {A review and classification of approaches for dealing with
  uncertainty in multi-criteria decision analysis for healthcare decisions}.
\newblock Pharmacoeconomics. 2015;33(5):445-55.
\newblock \href {http://dx.doi.org/10.1007/s40273-014-0251-x}
  {doi:10.1007/s40273-014-0251-x}.

\bibitem{Hendrickx2024}
{Hendrickx} K, {Perini} L, {Van der Plas} D, {Meert} W, {Davis} J.
\newblock {Machine learning with a reject option: a survey}.
\newblock Mach Learn. 2024;113(5):3073-110.
\newblock \href {http://dx.doi.org/10.1007/s10994-024-06534-x}
  {doi:10.1007/s10994-024-06534-x}.

\bibitem{Hu2023}
{Hu} XC, {Dai} M, {Sun} JM, {Sunderland} EM.
\newblock {The utility of machine learning models for predicting chemical
  contaminants in drinking water: promise, challenges, and opportunities}.
\newblock Curr Environ Health Rep. 2023;10(1):45-60.
\newblock \href {http://dx.doi.org/10.1007/s40572-022-00389-x}
  {doi:10.1007/s40572-022-00389-x}.

\bibitem{Geifman2017}
{Geifman} Y, {El-Yaniv} R.
\newblock {Selective classification for deep neural networks}.
\newblock In: {Guyon} I, {Von Luxburg} U, {Bengio} S, {Wallach} H, {Fergus} R,
  {Vishwanathan} S, et~al., editors. Advances in Neural Information Processing
  Systems 30 (NIPS 2017). Curran Associates; 2017. Available from:
  \url{https://proceedings.neurips.cc/paper_files/paper/2017/file/4a8423d5e91fda00bb7e46540e2b0cf1-Paper.pdf}
  [cited 2026 Sep 16].

\bibitem{GEMStat}
{UNEP GEMS/Water Programme}.
\newblock {GEMStat database} [Internet]; [cited 2026 Sep 1].
\newblock Available from: \url{https://gemstat.org/data-gemstat/}.

\bibitem{OWD-grdwtrsmpkwale}
{Loos} S, {Zhong} M, {Hope} R. {grdwtrsmpkwale: Groundwater analysis from 2016
  in Kwale, Kenya} [dataset]. Zenodo; 2023.
\newblock Version 0.0.1.
\newblock \href {http://dx.doi.org/10.5281/zenodo.10016875}
  {doi:10.5281/zenodo.10016875}.

\bibitem{OWD-boreholelabdata}
{Mhango} E, {Paterson} F, {Rattray} J. {boreholelabdata: Water Quality
  Laboratory Analysis and Reporting Dataset -- Malawi (2017--2019)} [dataset].
  Zenodo; 2025.
\newblock \href {http://dx.doi.org/10.5281/zenodo.17433461}
  {doi:10.5281/zenodo.17433461}.

\bibitem{OWD-mwgroundwaterdata}
{Mhango} E. {mwgroundwaterdata: Malawi Groundwater Monitoring Time-Series
  Dataset (2024--2025)} [dataset]. Zenodo; 2026.
\newblock \href {http://dx.doi.org/10.5281/zenodo.18798409}
  {doi:10.5281/zenodo.18798409}.

\bibitem{CGWB-NWDP}
{Central Ground Water Board}.
\newblock {Ground Water Quality (Manual - Chemical Parameters), Central Ground
  Water Board (CGWB)} [dataset]. National Water Data Portal, National Water
  Informatics Centre, Ministry of Jal Shakti; [cited 2026 Aug 7].
\newblock Available from:
  \url{https://nwdp.nwic.gov.in/dataset/7bb1e7c7-bcc1-48bb-8bcd-32c19473804a}.

\bibitem{CGWB-GWQ2019}
{Central Ground Water Board}.
\newblock {Ground Water Quality Data -- 2019} [data tables, PDF]. Ministry of
  Jal Shakti; 2019 [cited 2026 Sep 12].
\newblock Available from:
  \url{https://cgwb.gov.in/cgwbpnm/public/uploads/documents/16860545111615830419file.pdf}.
\newblock Internet Archive snapshot 20240124163108.

\bibitem{CGWB-GWQ2021}
{Central Ground Water Board}.
\newblock {Ground Water Quality Data -- 2021} [data tables, PDF]. Ministry of
  Jal Shakti; 2021 [cited 2026 Sep 12].
\newblock Available from:
  \url{https://cgwb.gov.in/cgwbpnm/public/uploads/documents/1686055122486176238file.pdf}.
\newblock Internet Archive snapshot 20231015091851.

\bibitem{CGWB-GWQ2022}
{Central Ground Water Board}.
\newblock {Ground Water Quality Data -- 2022 (Ground Water Quality of
  Unconfined Aquifers)} [data tables, PDF]. Ministry of Jal Shakti; 2022 [cited
  2026 Sep 12].
\newblock Available from:
  \url{https://cgwb.gov.in/sites/default/files/inline-files/final_nhs-wq_pre_2022_compressed.pdf}.

\bibitem{CGWB-AGWQR2025}
{Central Ground Water Board}.
\newblock {Annual ground water quality report-2025}.
\newblock Faridabad: Central Ground Water Board, Ministry of Jal Shakti,
  Government of India; 2025 Nov.
\newblock Available from:
  \url{https://cgwb.gov.in/cgwbpnm/public/uploads/documents/1762854375262680475file.pdf}
  [cited 2026 Sep 16].

\bibitem{WQP}
{Water Quality Portal}. Washington (DC): National Water Quality Monitoring
  Council, United States Geological Survey (USGS), Environmental Protection
  Agency (EPA); 2021.
\newblock \href {http://dx.doi.org/10.5066/P9QRKUVJ} {doi:10.5066/P9QRKUVJ}.

\bibitem{WHO2022GDWQ}
{World Health Organization}.
\newblock {Guidelines for drinking-water quality: fourth edition incorporating
  the first and second addenda}.
\newblock Geneva: World Health Organization; 2022.
\newblock ISBN 978-92-4-004506-4.

\bibitem{USEPA-NPDWR}
{United States Environmental Protection Agency}.
\newblock {National Primary Drinking Water Regulations} [Internet]. Washington
  (DC): US EPA; [cited 2026 Sep 15].
\newblock Available from:
  \url{https://www.epa.gov/ground-water-and-drinking-water/national-primary-drinking-water-regulations}.

\bibitem{Brown2001}
{Brown} LD, {Cai} TT, {DasGupta} A.
\newblock {Interval estimation for a binomial proportion}.
\newblock Stat Sci. 2001;16(2):101-33.
\newblock \href {http://dx.doi.org/10.1214/ss/1009213286}
  {doi:10.1214/ss/1009213286}.

\bibitem{Pedregosa2011}
{Pedregosa} F, {Varoquaux} G, {Gramfort} A, {Michel} V, {Thirion} B, {Grisel}
  O, et~al.
\newblock {Scikit-learn: machine learning in Python}.
\newblock J Mach Learn Res. 2011;12:2825-30.
\newblock Available from:
  \url{https://www.jmlr.org/papers/v12/pedregosa11a.html}.

\bibitem{Saito2015}
{Saito} T, {Rehmsmeier} M.
\newblock {The precision-recall plot is more informative than the ROC plot when
  evaluating binary classifiers on imbalanced datasets}.
\newblock PLoS One. 2015;10(3):e0118432.
\newblock \href {http://dx.doi.org/10.1371/journal.pone.0118432}
  {doi:10.1371/journal.pone.0118432}.

\bibitem{CGWB-ManualAquiferMapping}
{Central Ground Water Board}.
\newblock {Manual on aquifer mapping}.
\newblock Faridabad: Central Ground Water Board, Ministry of Water Resources;
  [date unknown].
\newblock Available from:
  \url{https://cgwb.gov.in/cgwbpnm/public/uploads/documents/16861227981023889265file.pdf}
  [cited 2026 Sep 15].

\bibitem{Poggio2021}
{Poggio} L, {de Sousa} LM, {Batjes} NH, {Heuvelink} GBM, {Kempen} B, {Ribeiro}
  E, et~al.
\newblock {SoilGrids 2.0: producing soil information for the globe with
  quantified spatial uncertainty}.
\newblock SOIL. 2021;7(1):217-40.
\newblock \href {http://dx.doi.org/10.5194/soil-7-217-2021}
  {doi:10.5194/soil-7-217-2021}.

\bibitem{Zanaga2022}
{Zanaga} D, {Van De Kerchove} R, {Daems} D, {De Keersmaecker} W, {Brockmann} C,
  {Kirches} G, et~al.
\newblock {ESA WorldCover 10 m 2021 v200 [dataset]}.
\newblock Zenodo; 2022.
\newblock \href {http://dx.doi.org/10.5281/zenodo.7254221}
  {doi:10.5281/zenodo.7254221}.

\bibitem{WorldPop2020}
{WorldPop}, {Bondarenko} M. {Individual countries 1km population density
  (2000-2020)} [dataset]. University of Southampton; 2020.
\newblock \href {http://dx.doi.org/10.5258/SOTON/WP00674}
  {doi:10.5258/SOTON/WP00674}.

\bibitem{Chakraborty2022}
{Chakraborty} S.
\newblock {TOPSIS and modified TOPSIS: a comparative analysis}.
\newblock Decis Anal J. 2022;2:100021.
\newblock \href {http://dx.doi.org/10.1016/j.dajour.2021.100021}
  {doi:10.1016/j.dajour.2021.100021}.

\bibitem{TervonenLahdelma2007}
{Tervonen} T, {Lahdelma} R.
\newblock {Implementing stochastic multicriteria acceptability analysis}.
\newblock Eur J Oper Res. 2007;178(2):500-13.
\newblock \href {http://dx.doi.org/10.1016/j.ejor.2005.12.037}
  {doi:10.1016/j.ejor.2005.12.037}.

\bibitem{BGS-Earthwise-Malawi}
{Upton} K, {{\'O} Dochartaigh} B{\'E}, {Chunga} B, {Bellwood-Howard} I.
\newblock {Africa Groundwater Atlas: hydrogeology of Malawi} [Internet].
  British Geological Survey; 2018 [updated 2023 Dec 20; cited 2026 Sep 15].
\newblock Available from:
  \url{https://earthwise.bgs.ac.uk/index.php/Hydrogeology_of_Malawi}.

\bibitem{George2012}
{George} CM, {Zheng} Y, {Graziano} JH, {Rasul} SB, {Hossain} Z, {Mey} JL,
  et~al.
\newblock {Evaluation of an arsenic test kit for rapid well screening in
  Bangladesh}.
\newblock Environ Sci Technol. 2012;46(20):11213-9.
\newblock \href {http://dx.doi.org/10.1021/es300253p} {doi:10.1021/es300253p}.

\bibitem{Reddy2020}
{Reddy} RR, {Rodriguez} GD, {Webster} TM, {Abedin} MJ, {Karim} MR, {Raskin} L,
  et~al.
\newblock {Evaluation of arsenic field test kits for drinking water:
  recommendations for improvement and implications for arsenic affected regions
  such as Bangladesh}.
\newblock Water Res. 2020;170:115325.
\newblock \href {http://dx.doi.org/10.1016/j.watres.2019.115325}
  {doi:10.1016/j.watres.2019.115325}.

\bibitem{Trimmer2024}
{Trimmer} JT, {Delaire} C, {Marshall} K, {Khush} R, {Peletz} R.
\newblock {Centralized or onsite testing? Examining the costs of water quality
  monitoring in rural Africa}.
\newblock Environ Sci Technol. 2024;58(26):11236-46.
\newblock \href {http://dx.doi.org/10.1021/acs.est.4c01916}
  {doi:10.1021/acs.est.4c01916}.

\bibitem{Pichel2023}
{Pichel} N, {Hymn{\^o} de Souza} F, {Sabogal-Paz} LP, {Shah} PK, {Adhikari} N,
  {Pandey} S, et~al.
\newblock {Field-testing solutions for drinking water quality monitoring in
  low- and middle-income regions and case studies from Latin American, African
  and Asian countries}.
\newblock J Environ Chem Eng. 2023;11(6):111180.
\newblock \href {http://dx.doi.org/10.1016/j.jece.2023.111180}
  {doi:10.1016/j.jece.2023.111180}.

\bibitem{Robertson2026-PLOSW0000593}
{Robertson} DJC, {Byrns} S, {Ellis} R, {White} CJ, {Nhlema} M, {Ngwira} G,
  et~al.
\newblock {Persistent groundwater data scarcity: challenges and advancements
  from a co-creation process in Malawi}.
\newblock PLOS Water. 2026;5(9):e0000593.
\newblock \href {http://dx.doi.org/10.1371/journal.pwat.0000593}
  {doi:10.1371/journal.pwat.0000593}.

\bibitem{Jameel2021}
{Jameel} Y, {Mozumder} MRH, {van Geen} A, {Harvey} CF.
\newblock {Well-switching to reduce arsenic exposure in Bangladesh: making the
  most of inaccurate field kit measurements}.
\newblock Geohealth. 2021;5(12):e2021GH000464.
\newblock \href {http://dx.doi.org/10.1029/2021GH000464}
  {doi:10.1029/2021GH000464}.

\bibitem{Ploton2020}
{Ploton} P, {Mortier} F, {R{\'e}jou-M{\'e}chain} M, {Barbier} N, {Picard} N,
  {Rossi} V, et~al.
\newblock {Spatial validation reveals poor predictive performance of
  large-scale ecological mapping models}.
\newblock Nat Commun. 2020;11(1):4540.
\newblock \href {http://dx.doi.org/10.1038/s41467-020-18321-y}
  {doi:10.1038/s41467-020-18321-y}.

\bibitem{SteyerbergVergouwe2014}
{Steyerberg} EW, {Vergouwe} Y.
\newblock {Towards better clinical prediction models: seven steps for
  development and an ABCD for validation}.
\newblock Eur Heart J. 2014;35(29):1925-31.
\newblock \href {http://dx.doi.org/10.1093/eurheartj/ehu207}
  {doi:10.1093/eurheartj/ehu207}.

\bibitem{VanCalster2016}
{Van Calster} B, {Nieboer} D, {Vergouwe} Y, {De Cock} B, {Pencina} MJ,
  {Steyerberg} EW.
\newblock {A calibration hierarchy for risk models was defined: from utopia to
  empirical data}.
\newblock J Clin Epidemiol. 2016;74:167-76.
\newblock \href {http://dx.doi.org/10.1016/j.jclinepi.2015.12.005}
  {doi:10.1016/j.jclinepi.2015.12.005}.

\bibitem{Purvis2026-PLOSW0000442}
{Purvis} T, {Nguyen} T, {McHugh} C, {Yeolekar} S, {Ding} E, {Wang} J, et~al.
\newblock {The Lishtot TestDrop Pro device cannot reliably indicate microbial
  water safety}.
\newblock PLOS Water. 2026;5(6):e0000442.
\newblock \href {http://dx.doi.org/10.1371/journal.pwat.0000442}
  {doi:10.1371/journal.pwat.0000442}.

\bibitem{Peletz2016}
{Peletz} R, {Kumpel} E, {Bonham} M, {Rahman} Z, {Khush} R.
\newblock {To what extent is drinking water tested in sub-Saharan Africa? A
  comparative analysis of regulated water quality monitoring}.
\newblock Int J Environ Res Public Health. 2016;13(3):275.
\newblock \href {http://dx.doi.org/10.3390/ijerph13030275}
  {doi:10.3390/ijerph13030275}.

\end{thebibliography}

\end{document}
